\documentclass[11pt,oneside]{book}
\usepackage[a4paper,margin=27mm,headheight=15pt]{geometry}
\usepackage[T1]{fontenc}
\usepackage[utf8]{inputenc}
\usepackage{lmodern,microtype}
\usepackage{amsmath,amssymb,amsthm,mathtools,bm}
\usepackage{graphicx,booktabs,longtable,array,tabularx}
\usepackage{enumitem,fancyhdr,xcolor,listings}
\usepackage{xurl}
\usepackage[unicode,colorlinks=true,linkcolor=blue!45!black,citecolor=blue!45!black,urlcolor=blue!45!black]{hyperref}
\usepackage{tikz}
\hypersetup{pdftitle={Conserving Probability in a Chaotic State Space},pdfsubject={Lorenz-63 Fokker-Planck solver: mathematical foundations and evidence-based development},pdfauthor={Prepared for Adam James Sheppard},pdfkeywords={Lorenz, Fokker-Planck, DG, SIPG, Crank-Nicolson, positivity, QP, MFEM, AMR}}
\usetikzlibrary{arrows.meta,positioning}
\graphicspath{{../../output/pdf/solver_monograph/}}
\setlength{\parskip}{5pt}
\setlength{\parindent}{1em}
\setlength{\emergencystretch}{3em}
\setlist{itemsep=3pt,topsep=4pt}
\pagestyle{fancy}\fancyhf{}\fancyhead[L]{\small Lorenz probability-density solver}\fancyhead[R]{\small Mathematical and computational record}\fancyfoot[C]{\thepage}
\renewcommand{\headrulewidth}{.3pt}
\newtheorem{theorem}{Theorem}[chapter]
\newtheorem{proposition}[theorem]{Proposition}
\newtheorem{lemma}[theorem]{Lemma}
\newtheorem{corollary}[theorem]{Corollary}
\theoremstyle{definition}\newtheorem{definition}[theorem]{Definition}\newtheorem{example}[theorem]{Example}
\theoremstyle{remark}\newtheorem{remark}[theorem]{Remark}
\newcommand{\R}{\mathbb R}\newcommand{\E}{\mathbb E}\newcommand{\Prob}{\mathbb P}
\newcommand{\dd}{\,\mathrm d}\newcommand{\grad}{\nabla}\newcommand{\divg}{\nabla\!\cdot}
\newcommand{\jump}[1]{[\![#1]\!]}\newcommand{\avg}[1]{\{#1\}}
\newcommand{\code}[1]{\nolinkurl{#1}}
\newcommand{\evidence}[1]{\par\smallskip\noindent\textbf{Project evidence.} #1\par}
\newcommand{\scope}[1]{\par\smallskip\noindent\textbf{Scope of the claim.} #1\par}
\lstset{basicstyle=\ttfamily\small,breaklines=true,columns=fullflexible,frame=single,framerule=.2pt,showstringspaces=false}
\begin{document}
\frontmatter
\begin{titlepage}
\hypersetup{pageanchor=false}
\centering\vspace*{25mm}
{\Huge\bfseries Conserving Probability\\[5mm] in a Chaotic State Space\par}
\vspace{12mm}{\Large A mathematical and computational monograph\\on the Lorenz--63 Fokker--Planck solver\par}
\vspace{18mm}{\large Foundations, discretisation, positivity, adaptive resolution,\\verification and the complete recorded development history\par}
\vfill
{\large Adam James Sheppard's research project\par}
\vspace{6mm}{\large Evidence snapshot: 3 October 2026\par}
\vspace{6mm}{\normalsize Solver revision \code{1781b40ad156f01702aeeb16eaf7e0240ede7378}\par}
\vspace{6mm}{\small Source repository: \url{https://github.com/AdamJamesSheppard/lorenz-training-set}\par}
\vspace{12mm}{\small Prepared from repository source, immutable local experiments, numerical arrays,\\version history and verified primary literature. No solver equations were changed.\par}
\end{titlepage}
\hypersetup{pageanchor=true}
\chapter*{Abstract}\addcontentsline{toc}{chapter}{Abstract}
This monograph develops the mathematics and implementation of a conservative three-dimensional probability-density solver for stochastic Lorenz--63 dynamics. It starts with probability measures and stochastic state evolution, derives the Kolmogorov forward equation, and follows its flux structure into discontinuous Galerkin discretisation. Conservative upwind transport and full-tensor symmetric interior-penalty diffusion lead to a nonsymmetric semidiscrete system. Crank--Nicolson provides an accurate high-order target, while a global nonnegative-average repair and a cell-local mass-matrix quadratic programme restore an admissible probability density.

The development record supplies the causal structure of the book. A confounded equal-DOF comparison was replaced by same-mesh tests. A roundoff repair violated a limiter invariant and was corrected. Global scaling degraded an accurate Q2 trajectory; local QP correction performed better. Startup temporal and full-SPD spatial gates passed, but mature bimodal densities exposed genuine within-cell negativity and under-resolution. Certificate-depth, average-level AFC and Bernstein-DOF diagnostics narrowed the mechanism. Uniform and graded refinements showed that resolution controls the intervention. A basis-transformed MFEM port reproduced DOLFINx before static nonconforming hexahedral refinement was enabled. Subsequent support, depth, timestep and domain comparisons established bounded sensitivity for one mature posterior and a short forecast.

The final tested implementation is a static NC-hex Q2--CN/full-SPD/local-QP solver. Both aligned domain expansions pass; the second-domain common-interior voxel difference independently recomputes to $1.9549638427469373\times10^{-12}$. This is evidence of boundary insensitivity for the tested law and horizon, not a continuum error bound. The book distinguishes exact-arithmetic algebra, sufficient polynomial certificates, floating-point acceptance criteria, reference uncertainty and empirical convergence. It includes a complete specification, pseudocode, code-to-mathematics mapping, reproducibility instructions, a failed-method catalogue and a hashed source ledger. Large-scale dataset certification remains withheld; the recent discussion supports a bounded provisional neural-operator pilot without converting that judgement into a convergence theorem.

\chapter*{How to read this book}\addcontentsline{toc}{chapter}{How to read this book}
The reader is assumed to know calculus, linear algebra and elementary probability. Specialist concepts are developed when the solver first needs them. A result labelled \emph{Theorem} or \emph{Proposition} is proved here under stated assumptions. A numerical result is identified as project evidence. A heuristic diagnosis is kept separate from the quantity actually measured.

The record begins with the initial Git baseline on 9 September 2026. Files present in that baseline may predate it; their existence alone cannot establish an earlier sequence of experiments. The available conversation supplies context and recent user decisions, but repository files and immutable numerical records take precedence for numerical values. The book does not invent a pre-Git development log. No confirmed Chang--Cooper production branch appears in the current source tree: that family is explained as a literature comparator, while the implemented finite-volume reference uses upwind drift, diagonal diffusion and SSPRK time stepping.

Each part returns to the same question: what does this particular piece of mathematics require from the algorithm? The answer becomes a code contract, then an experimental test. A passing test only establishes what the test actually measures. Conversely, a failed aggregate gate may contain valuable successful components. This distinction is essential for interpreting the mature-state and mesh-grading episodes.

The evidence ledger indexes 883 compact source/experiment records and 86 distinct local run directories at the audited snapshot, with 69 recorded commits. Those counts refer to the evidence snapshot before monograph creation. Large arrays were selectively audited; the inventory is comprehensive for tracked paths and available compact records, not an assertion that every byte of every binary archive was interpreted. The figures drawn from saved density arrays are conservative voxel marginals, not invented attractor densities. Conceptual diagrams are explicitly labelled as such.

\tableofcontents\listoffigures\listoftables
\chapter*{Notation and evidence conventions}\addcontentsline{toc}{chapter}{Notation and evidence conventions}
\begin{longtable}{p{.20\textwidth}p{.73\textwidth}}
\toprule Symbol & Meaning\\\midrule\endhead
$X_t\in\R^3$ & Random Lorenz state at time $t$.\\
$x=(x_1,x_2,x_3)$ & A deterministic point in state space; also written $(x,y,z)$ where convenient.\\
$p(x,t)$ & Density with respect to Lebesgue volume, not probability per voxel.\\
$f(x)$ & Lorenz drift; $\sigma=10$, $\rho=28$, $\beta=8/3$.\\
$B,D$ & Noise factor and diffusion tensor, $D=BB^T/2$.\\
$\Omega$ & Finite state-space box; default $[-30,30]\times[-40,40]\times[-10,70]$.\\
$J$ & Total probability flux $fp-D\grad p$.\\
$K,F,n$ & A mesh cell, interior face and an outward unit normal.\\
$Q_k$ & Tensor polynomials of degree at most $k$ in each coordinate.\\
$V_h$ & Broken DG finite-element space.\\
$M,S,K_h$ & Mass matrix, positive-sign spatial form matrix, and evolution matrix $K_h=-S$.\\
$c,u$ & Local polynomial and global coefficient vectors.\\
$H,a,C$ & Reference local mass metric, average functional and Bernstein constraint matrix.\\
$w_K,x_K$ & Raw and repaired cell averages.\\
$\Delta t,\tau$ & Numerical timestep and forecast horizon.\\
$P_Vp$ & Piecewise-constant field of conservative voxel averages.\\
$\mathrm{TV}(p,q)$ & $\frac12\|p-q\|_{L^1}$ for normalized densities on the same space.\\
$C_n$ & Stepwise relative raw-to-corrected $L^1$ change.\\
\bottomrule\end{longtable}
\noindent\textbf{Evidence labels.} \emph{Established} identifies a derivation or an imported theorem with hypotheses. \emph{Formal} identifies a discretisation property in exact arithmetic. \emph{Numerical} identifies measurements at declared tolerances. \emph{Heuristic} identifies an explanation supported but not implied by the data. \emph{Open} identifies a missing proof, reference, test or scope extension. Repository assumption identifiers are discussed in Appendix~\ref{app:assumptions}; they are not a substitute for locally stating the relevant hypothesis.
\mainmatter
\part{From uncertain states to a conservative PDE}
\chapter{The problem the solver actually solves}\label{ch:purpose}
\section{Why an individual trajectory is insufficient}
Lorenz--63 supplies three coupled state coordinates, yet a state estimator rarely knows all three exactly. An observation may constrain $z$ while leaving two lobes in $(x,y)$ plausible. A covariance matrix describes spread around a mean, but a mean near the gap between lobes can correspond to very little probability. The object required for propagation is therefore a probability law, with enough structure to represent multimodality and tails.

The solver computes how that law changes under a specified stochastic differential equation. Its output is a function over the full three-dimensional state space. A point in that space is a possible system state; a computational cell is a small collection of possible states. It is not a parcel of physical atmospheric fluid. This distinction gives the boundary condition its probabilistic interpretation: boundary flux measures the probability assigned to states crossing the artificial computational box.

There are two broad computational representations. An ensemble advances many sample paths and estimates probabilities from the samples. A density solver advances a deterministic function satisfying a forward equation. The ensemble retains path information and avoids a full state-space mesh; the density solver resolves a field directly and supports a likelihood update without reducing it to moments. Neither representation eliminates numerical error. Sampling error, ensemble integration bias and density-reconstruction bias enter the first; spatial, temporal, boundary, projection and algebraic errors enter the second.

The project uses both. DG density forecasts are the candidate training targets, while Monte Carlo particles supply independent moment and distribution comparisons. The comparison must use the same initial law. Drawing particles from an analytic Gaussian while evolving an inaccurately represented finite-element approximation confounds initial representation with dynamical error. This issue was consequential in the original method ranking and motivates the exact-native-law samplers discussed later.

\section{A forecast within a Bayesian cycle}
Let $p_{k|k}$ denote the density after observations through time $t_k$. Propagation over $\tau=t_{k+1}-t_k$ produces $p_{k+1|k}$. An observation $y_{k+1}$ then gives
\begin{equation}\label{eq:bayes-overview}
p_{k+1|k+1}(x)=\frac{\ell(y_{k+1}|x)p_{k+1|k}(x)}{\int_\Omega \ell(y_{k+1}|s)p_{k+1|k}(s)\dd s}.
\end{equation}
The training contract is the map from the complete posterior to the complete forecast. It is not a map from a posterior mean to a forecast mean. The early dataset scaffold stores posterior and forecast arrays, records failed and burn-in cycles, and splits independent trajectories rather than adjacent timesteps. Those engineering contracts remain useful even though the mature MFEM implementation is not wired into that serial DOLFINx writer as a production generator.

\section{The causal architecture}
The mathematical path is short enough to state before deriving it. Stochastic increments specify a generator. Its adjoint governs density. Density is a conserved nonnegative quantity. Conservation motivates flux-based element coupling. Local polynomial approximation improves resolution but can generate negative values. Implicit integration controls stiffness without automatically controlling positivity. Constrained repair restores admissibility, and the magnitude of that repair becomes evidence about the underlying approximation. Mature laws require local resolution, prompting mesh adaptivity rather than indefinitely stronger limiting.

\begin{figure}[ht]\centering
\begin{tikzpicture}[node distance=6mm, every node/.style={draw,rounded corners,align=center,font=\small,text width=10.5cm},>=Stealth]
\node(a){Uncertain Lorenz state: a probability law};
\node(b)[below=of a]{Stochastic increments $\longrightarrow$ generator $\longrightarrow$ forward PDE};
\node(c)[below=of b]{Total probability flux $\longrightarrow$ conservative DG interfaces};
\node(d)[below=of c]{Q2 and Crank--Nicolson $\longrightarrow$ accurate raw target with possible negativity};
\node(e)[below=of d]{Average repair and local QP $\longrightarrow$ admissible density};
\node(f)[below=of e]{Correction and density differences $\longrightarrow$ resolution diagnosis and static AMR};
\draw[->](a)--(b);\draw[->](b)--(c);\draw[->](c)--(d);\draw[->](d)--(e);\draw[->](e)--(f);
\end{tikzpicture}
\caption{Conceptual dependency diagram. It illustrates mathematical causality and contains no simulated numerical result.}\label{fig:causality}
\end{figure}

\evidence{\code{lorenz_fpe/core.py}, \code{lorenz_fpe/dataset.py}, \code{mfem/mfem_physical_equivalence.cpp} and \code{sample_dataset/dataset_schema.json} establish the implemented state and output contracts.}

\chapter{Probability as a computational invariant}\label{ch:probability}
\section{Measures, random vectors and densities}
A probability space $(\mathcal S,\mathcal F,\Prob)$ consists of outcomes, measurable events and a countably additive measure of total mass one. A random vector $X:\mathcal S\to\R^3$ is measurable. Its law is the pushforward measure $\mu(A)=\Prob(X\in A)$ on Borel sets $A\subset\R^3$. The outcomes are never enumerated by the density solver; it operates on the law.

If $\mu$ is absolutely continuous with respect to Lebesgue measure, there exists an integrable density $p\ge0$ satisfying
\begin{equation}\label{eq:density-def}
\mu(A)=\int_A p(x)\dd x,\qquad \int_{\R^3}p(x)\dd x=1.
\end{equation}
Absolute continuity is an assumption, not a property of every random vector. A point mass has no ordinary Lebesgue density. Nondegenerate diffusion provides smoothing for positive times under suitable conditions, but a rank-deficient noise matrix does not automatically provide the same elliptic argument. The implemented constant full-SPD experiment avoids that degeneracy.

A density has units inverse to state-space volume. If a voxel has volume $V$, its probability is $V\bar p$, not $\bar p$. The distinction becomes crucial when refining a mesh or comparing nonuniform cells: adding density values without volume weights is not a probability integral. This is why the global average-repair problem uses actual cell volumes on AMR meshes.

\section{Expectations and moments}
For measurable $g$ with $\int|g|p<\infty$, define $\E[g(X)]=\int gp$. The mean $\mu\in\R^3$ and covariance $\Sigma\in\R^{3\times3}$ are
\begin{equation}\label{eq:moments}
\mu_i=\int x_i p(x)\dd x,\quad
\Sigma_{ij}=\int(x_i-\mu_i)(x_j-\mu_j)p(x)\dd x.
\end{equation}
On the finite box these integrals exist for every integrable density, since the coordinate functions are bounded. Covariance is symmetric positive semidefinite: for every vector $a$, $a^T\Sigma a=\E[(a^T(X-\mu))^2]\ge0$. A materially negative covariance eigenvalue therefore indicates a defective statistical calculation, a signed distribution, or numerical roundoff beyond the intended tolerance.

Means and covariances are summaries. Consider a one-dimensional mixture of equal narrow Gaussians centred at $\pm a$. Its mean is zero and its variance is $a^2+s^2$. A single Gaussian with that variance shares the first two moments while assigning entirely different probability near the centre. In three dimensions the Lorenz lobe geometry makes this discrepancy operational: a posterior can remain ambiguous across lobes even after a $z$ observation.

\section{Marginals and conditioning}
The $x$ marginal is $p_x(x)=\int p(x,y,z)\dd y\dd z$. Tonelli's theorem justifies exchanging these integrals for nonnegative $p$. For integrable signed numerical fields, Fubini requires absolute integrability. Marginalisation preserves total mass and decreases $L^1$ discrepancy:
\begin{equation}\label{eq:marginal-contraction}
\int\left|\int(p-q)\dd y\dd z\right|\dd x\le\int|p-q|\dd x\dd y\dd z.
\end{equation}
Thus a small full-density error controls marginal errors; small marginal errors cannot control the full joint dependence. Two densities can share every one-dimensional marginal and have different correlations or multimodal connections.

Conditioning on a continuous observation is expressed through a likelihood, not by restricting to an event of positive probability $X=x$. For Gaussian measurement noise, $\ell(y|x)$ is a positive function. Multiplication by that likelihood changes the law and normalization restores total probability. In a discrete solver the product must be projected into the finite-element space. Projection can introduce oscillations even though the analytic product is positive; positivity protection is consequently needed after analysis as well as propagation.

\section{Signed fields and the meaning of negative mass}
Let $p^- =\max(-p,0)$ and $p^+=\max(p,0)$. A signed field can integrate to one while containing a negative region because
\begin{equation}
\int p=\int p^+-\int p^-=1.
\end{equation}
The negative mass $m^-=\int p^-$ is an admissibility defect. It is not removed by dividing by $\int p$. Pointwise clipping removes it but increases total mass; renormalizing the clipped field changes all cells. The project instead repairs averages and polynomial shape with explicit conservation constraints.

A quadrature estimate of $m^-$ can miss a narrow negative region. A sufficient whole-cell polynomial certificate addresses that gap. The recorded phrase ``zero measured negative mass'' is therefore accompanied by certification counts and tolerance information rather than presented alone as a theorem.

\begin{example}[Conservation without admissibility]
On two unit-volume cells, averages $(1.2,-0.2)$ have mass one. The nonnegative equal-mass Euclidean projection is $(1,0)$. The total mass was already correct; positivity required a redistribution. Within a single cell, a positive average can coexist with a negative quadratic polynomial near an interior point. The two defects require different constraints.
\end{example}

\chapter{Lorenz dynamics and the origin of difficult densities}\label{ch:lorenz}
\section{Equations and equilibria}
The deterministic drift is
\begin{equation}\label{eq:lorenz}
f(x,y,z)=\begin{pmatrix}\sigma(y-x)\\x(\rho-z)-y\\xy-\beta z\end{pmatrix},
\qquad (\sigma,\rho,\beta)=(10,28,8/3).
\end{equation}
The coordinates are nondimensional in this project. The original Lorenz model and its sensitivity motivation are established in Lorenz's paper~\cite{lorenz}. We use the equations as a nonlinear state-estimation benchmark, rather than claiming quantitative atmospheric validation.

The origin is an equilibrium. For $\rho>1$, two further equilibria are
\begin{equation}
E_\pm=(\pm\sqrt{\beta(\rho-1)},\ \pm\sqrt{\beta(\rho-1)},\ \rho-1).
\end{equation}
To derive them, the first equation gives $y=x$. The third gives $z=x^2/\beta$. The second becomes $x(\rho-1-x^2/\beta)=0$. For the chosen parameters the nonzero equilibria lie at $(\pm\sqrt{72},\pm\sqrt{72},27)$. This geometry motivates symmetric lobe components, but does not say that a mature stochastic density equals two Gaussians at these points.

The drift has symmetry $f(Rx)=Rf(x)$ for $R=\operatorname{diag}(-1,-1,1)$. A symmetric initial law remains symmetric in the deterministic model. For additive noise, symmetry of the generator additionally requires $RDR^T=D$. The tested full tensor has nonzero $D_{13}$ and $D_{23}$, so that condition fails. Symmetrizing the mature mixture is a chosen initial-law construction; it does not establish that the full-SPD stochastic invariant law has exact lobe symmetry.

\section{Compressibility and dissipation}
The divergence is constant:
\begin{equation}\label{eq:divergence}
\divg f=-\sigma-1-\beta=-41/3.
\end{equation}
For the deterministic flow $\Phi_t$ where it exists smoothly, Jacobi's formula gives
$\frac{d}{dt}\det D\Phi_t=(\divg f)(\Phi_t)\det D\Phi_t$.
Hence local state-space volume contracts as $e^{-41t/3}$. This concerns volume, not probability: probability is transported into a shrinking volume and density can increase.

To see that boundedness is compatible with nonlinear coupling, take
$V=\rho x^2+\sigma y^2+\sigma(z-2\rho)^2$. Differentiating along~\eqref{eq:lorenz} cancels the $xyz$ and $xy$ cross terms and yields
\begin{equation}\label{eq:lyapunov}
\dot V=-2\sigma\rho x^2-2\sigma y^2-2\sigma\beta(z-\rho)^2+2\sigma\beta\rho^2.
\end{equation}
Since $(z-\rho)^2\ge\frac12(z-2\rho)^2-\rho^2$, there exist positive $c,C$ depending on the parameters with $\dot V\le-cV+C$. The scalar comparison equation then bounds $V(t)$ by $e^{-ct}V(0)+C/c$. This is an absorbing-set argument for the deterministic system, not a proof of a particular chaotic invariant measure.

For additive noise, It\^o's formula introduces $\operatorname{tr}(D\nabla^2V)$, a constant. A stopped-process Lyapunov argument gives the same type of expectation bound and helps exclude finite-time explosion. Local Lipschitz continuity of the polynomial drift supplies local existence/uniqueness; the Lyapunov bound is needed to extend the solution globally. Merely calling the drift smooth would omit the growth issue.

\section{Stretching, folding and resolution}
Negative divergence does not mean contraction in every direction. The Jacobian of the drift contains off-diagonal couplings. Perturbations satisfy $\dot\delta= Df(X_t)\delta$ to first order, and directions can stretch while total volumes shrink. This produces narrow extended probability structures. Noise broadens those structures, but on a short interval anisotropic drift and mixed diffusion still generate strongly non-Gaussian shapes.

The deterministic attractor is a long-time geometric object. A stochastic density with nondegenerate noise is a volume distribution, not a set of zero volume. A short forecast from a conditioned Gaussian mixture need not display the full familiar butterfly outline. The saved initial density represents the chosen posterior law, constructed from a spun-up ensemble and fitted components; it should not be judged by whether an isosurface resembles an artistic attractor rendering.

\evidence{The mature configuration specifies 200,000 spin-up particles, spin-up time 2, spin-up step 0.001, eight components per lobe and a $z=25$ likelihood with variance 16. These are experiment choices from \code{experiments/mature-state-decision.yaml}; they do not establish stationary sampling or exactness of the mixture fit.}

\chapter{Stochastic increments and the diffusion convention}\label{ch:sde}
\section{Brownian motion and the integral equation}
A standard $m$-dimensional Wiener process $W_t$ has independent increments, continuous paths, $W_0=0$ and $W_{t+h}-W_t\sim N(0,hI_m)$. Its pathwise derivative does not exist as an ordinary smooth function. The SDE notation
\begin{equation}\label{eq:sde}
\mathrm dX_t=f(X_t)\dd t+B\mathrm dW_t
\end{equation}
means $X_t=X_0+\int_0^tf(X_s)\dd s+\int_0^tB\mathrm dW_s$. For constant $B$, the second integral is $BW_t$. The adaptedness and integrability conditions defining an It\^o integral are part of the stochastic model; a Gaussian random number inserted without the correct $\sqrt h$ scaling would define a different process.

For a short step conditional on $X_t=x$, the leading approximation is
\begin{equation}
X_{t+h}-x=f(x)h+B\sqrt h\,\xi+\text{higher-order terms},\quad \xi\sim N(0,I_m).
\end{equation}
The leading conditional covariance is $hBB^T$. The forward PDE convention uses half of this covariance rate:
\begin{equation}\label{eq:diffusion-convention}
D=\tfrac12BB^T.
\end{equation}
This factor is a recurring audit point. In the original identity-noise experiments $B=I$ implies $D=I/2$. In the later mixed-diffusion experiment the diagonal entries of $D$ are one. It is not the same noise magnitude with additional off-diagonal entries.

\section{From a target tensor to a noise factor}
For a real symmetric positive-definite $D$, Cholesky factorization gives $D=LL^T$ with positive diagonal entries of $L$. Taking $B=\sqrt2L$ realizes~\eqref{eq:diffusion-convention}. The tested tensor is
\begin{equation}\label{eq:spd-tensor}
D=\begin{pmatrix}1&0.4&0.2\\0.4&1&0.3\\0.2&0.3&1\end{pmatrix}.
\end{equation}
Its eigenvalues are approximately $(0.581212,0.811321,1.607467)$. The smallest is strictly positive, so $v^TDv\ge\lambda_{\min}\|v\|^2$. Its noise factor is approximately
\begin{equation}
B=\begin{pmatrix}1.4142135624&0&0\\0.5656854249&1.2961481397&0\\0.2828427125&0.3394673699&1.3434142715\end{pmatrix}.
\end{equation}
Using a square root other than Cholesky gives the same law for constant additive Gaussian noise when $BB^T$ agrees. The implementation stores $B$ in the Python model and derives $D$; the MFEM comparator uses the fixed tensor directly. Their equivalence checks therefore include the physical diffusion convention.

Mixed terms matter. For constant symmetric $D$, $\divg(D\grad p)=\sum_iD_{ii}p_{x_ix_i}+2\sum_{i<j}D_{ij}p_{x_ix_j}$. Omitting the factor two or using only diagonal entries changes the operator. A covariance check limited to diagonal entries can miss such a change; off-diagonal covariance and correlations were explicitly included in the full-SPD gate.

\section{It\^o's formula and the generator}
For $\varphi\in C^2$ with appropriate growth/integrability, It\^o's formula gives
\begin{equation}\label{eq:ito}
\mathrm d\varphi(X_t)=\grad\varphi(X_t)^Tf(X_t)\dd t
+\tfrac12\operatorname{tr}(BB^T\nabla^2\varphi(X_t))\dd t
+\grad\varphi(X_t)^TB\mathrm dW_t.
\end{equation}
The martingale term has zero expectation when square-integrable. The generator acting on observables is
\begin{equation}\label{eq:generator}
L\varphi=f\cdot\grad\varphi+D:\nabla^2\varphi.
\end{equation}
Here $A:C=\sum_{ij}A_{ij}C_{ij}$. The conditional expectation definition is
\begin{equation}
 L\varphi(x)=\lim_{h\downarrow0}\frac{\E[\varphi(X_h)|X_0=x]-\varphi(x)}h.
\end{equation}
It agrees with~\eqref{eq:generator} under the stated regularity. The stochastic integral and generator conventions are developed in the verified SDE reference~\cite{sarkka}; the subsequent derivation is given explicitly rather than delegated to that citation.

\section{Monte Carlo propagation and its errors}
Euler--Maruyama takes $X_{n+1}=X_n+f(X_n)h+B\sqrt h\xi_n$. Particle counts reduce statistical uncertainty but do not remove timestep bias. Common random numbers across comparisons reduce the variance of differences by correlating the approximations. A bootstrap estimates sampling variability under its resampling model; increasing bootstrap replicates improves estimation of that variability, while increasing particle count can reduce it. Neither operation proves that ensemble integration or density smoothing is unbiased.

The strengthened project reference uses up to one million paths and fixed bootstrap seeds. This was necessary because a finest deterministic covariance discrepancy of 0.00234 lay below a 200,000-path bootstrap p95 of 0.00649. Comparing numbers below the reference uncertainty floor cannot demonstrate further solver improvement. Deterministic density differences remain useful when covariance resolution is exhausted.

\chapter{Deriving the forward equation}\label{ch:forward}
\section{Weak evolution before a classical PDE}
Take a compactly supported smooth $\varphi$. Under nonexplosion and integrability, expectation of~\eqref{eq:ito} yields
\begin{equation}\label{eq:weak-law}
\frac{\mathrm d}{\mathrm dt}\int\varphi(x)p(x,t)\dd x
=\int\big(f\cdot\grad\varphi+D:\nabla^2\varphi\big)p\dd x.
\end{equation}
This can first be interpreted in integrated time form. Writing a time derivative under the integral needs enough time regularity, or a distributional interpretation. A law may be defined even before a smooth density is justified; weak evolution is therefore the appropriate starting point.

For constant $D$, integrate once by parts in the drift and twice in diffusion. Compact support removes boundary terms on $\R^3$:
\begin{align}
\int f_i p\,\partial_i\varphi\dd x&=-\int\varphi\,\partial_i(f_i p)\dd x,\\
\int D_{ij}p\,\partial_i\partial_j\varphi\dd x&=\int\varphi\,D_{ij}\partial_i\partial_jp\dd x.
\end{align}
The adjoint of $L$ acting on densities is consequently
\begin{equation}\label{eq:forward}
\partial_t p=L^*p=-\divg(fp)+\divg(D\grad p).
\end{equation}
For spatially varying noise, the adjoint is $\frac12\partial_i\partial_j((BB^T)_{ij}p)$ rather than automatically $\divg(D\grad p)$. The project's constant matrix assumption is essential to its current formula. Changing $B$ to depend on state requires a fresh model and flux audit.

\section{Compressibility cannot be omitted}
Expanding the drift gives $-\divg(fp)=-f\cdot\grad p-(\divg f)p$. The Lorenz divergence~\eqref{eq:divergence} is nonzero. A code evolving $p_t=-f\cdot\grad p+\divg(D\grad p)$ would transport a scalar concentration along characteristics while omitting the compressibility factor. It would not evolve the intended probability density.

In the diffusion-free smooth case, the characteristic identity is
\begin{equation}
p(\Phi_t(x),t)\det D\Phi_t(x)=p(x,0).
\end{equation}
This follows by differentiating both factors along the flow. The density grows where volume contracts so that probability stays fixed. Conservative weak advection incorporates this automatically through $-\int fp\cdot\grad v$ and interface fluxes; an additional explicit $-(\divg f)p$ term there would double count it.

\section{Moment evolution and unclosed statistics}
On whole space with sufficient decay and finite moments, testing with $x_i$ gives $\dot\mu_i=\E[f_i(X)]$. Testing with $x_ix_j$ gives
\begin{equation}
\frac{\mathrm d}{\mathrm dt}\E[X_iX_j]=\E[X_if_j+X_jf_i]+2D_{ij}.
\end{equation}
Because Lorenz drift is quadratic, these expressions involve higher moments. They do not form a closed Gaussian system. This explains why advancing mean and covariance alone cannot reproduce the full forecast without an additional approximation.

On a reflecting finite box, integration by parts for nonconstant test functions leaves boundary contributions. The mass test is special because its gradient vanishes. Whole-space moment identities must not be copied unchanged into a bounded-domain verifier without checking boundary terms. In the tested case negligible boundary probability supports approximate agreement, but it does not change the mathematical boundary model.

\section{Regularity and positivity}
The principal part is uniformly elliptic when $D$ is SPD. On a bounded box the drift is smooth and bounded; diffusion has positive energy $\int\grad p^TD\grad p$. Classical maximum-principle arguments, or positivity of the corresponding reflected diffusion semigroup, give nonnegative evolution under their boundary and regularity hypotheses. Corners of a box require care for classical boundary regularity; weak solutions provide a natural framework.

This continuous positivity property does not transfer automatically to an arbitrary Galerkin matrix. Basis functions change sign, the consistent mass inverse couples coefficients, and CN can have negative amplification for stiff modes. The rest of the book derives how conservation is retained discretely and how positivity is imposed separately.

\chapter{Flux, boundaries and what domain tests establish}\label{ch:flux}
\section{Local and global balance}
Define $J=fp-D\grad p$. Then $p_t+\divg J=0$. Integrating over a fixed region $A$ and applying the divergence theorem gives
\begin{equation}\label{eq:balance}
\frac{\mathrm d}{\mathrm dt}\int_Ap\dd x=-\int_{\partial A}J\cdot n\dd s.
\end{equation}
The sign is determined by the outward normal. Positive $J\cdot n$ removes probability from $A$. Internal face contributions cancel when regions are combined, provided the same numerical flux is used with opposite normal orientation on the two sides.

For $A=\Omega$, imposing $J\cdot n=0$ gives mass conservation. This is a condition on the \emph{total} flux. If drift pushes towards a wall, a diffusive gradient can balance it. Imposing separately $(fp)\cdot n=0$ and $(D\grad p)\cdot n=0$ would in general overconstrain or change the model. The implemented weak form omits exterior total-flux contributions; it does not independently prescribe two zero boundary fluxes.

\section{A one-dimensional boundary example}
For constant drift $a$ and diffusion $d>0$ on $[0,L]$, zero stationary flux satisfies $ap-dp'=0$. Thus $p(x)=Ce^{ax/d}$ with $C$ chosen to normalize. If $a\ne0$, the gradient at the boundary is nonzero. A homogeneous diffusion Neumann condition $p'=0$ would be incompatible with that zero-total-flux equilibrium. The example isolates why the probabilistic boundary convention matters before the full Lorenz tensor is introduced.

\section{The finite box changes the stochastic problem}
The original whole-space SDE permits paths outside the computational box. A reflecting FPE corresponds to a bounded-domain diffusion with an appropriate reflection mechanism; for anisotropic diffusion the associated conormal boundary structure must be respected. The solver is consequently a truncated approximation to the whole-space law, not a literal implementation of unconstrained trajectories on all $\R^3$.

Boundary placement error is small when the law has negligible interaction with the artificial boundary over the tested horizon. A boundary-layer probability is a useful indicator, but small mass in one layer does not by itself bound every observable. Domain expansion compares the solution while preserving physical interior resolution, isolating sensitivity to boundary placement.

\section{The completed aligned expansions}
The base box has volume $60\cdot80\cdot80=384000$. The broader AMR mesh has background dimensions $45\times54\times54$. Adding $r$ background cells on each side produces extents
\begin{equation}
[-30-r(60/45),30+r(60/45)]\times[-40-r(80/54),40+r(80/54)]\times[-10-r(80/54),70+r(80/54)].
\end{equation}
The two tested values are $r=1,2$. Background counts become $47\times56\times56$ and $49\times58\times58$. Original interior coordinates and refinement marks are retained. The export adds four voxels per background cell per side, so shapes are $188\times224\times224$ and $196\times232\times232$.

For the second comparison, removing four outer voxels per axis produces precisely the first-expanded grid. Independent recomputation gives $1.9549638427469373\times10^{-12}$ for the volume-weighted $L^1$ difference and zero directly summed mass in the added shell. The report's $6.66\times10^{-16}$ outside-box probability comes from subtracting nearly equal totals and is a roundoff residual. These statements agree within arithmetic precision, but the direct shell sum is the clearer measurement.

\scope{The result establishes numerical boundary insensitivity for this mature law, this $D$ and $t=0.05$. It does not establish spatial convergence, long-time accuracy, or independence for new posteriors. Both expanded runs retain the same interior discretisation and can share the same interior bias.}
\evidence{\code{runs/mfem-domain-sensitivity/20261003T143753Z/domain_decision.json}; reused first expansion \code{20261003T113423Z/padding1}; source ledger records both final-array hashes and the independent recomputation.}

\chapter{Bayesian laws and the mature initial condition}\label{ch:bayes}
\section{Likelihood and evidence}
For $y=Hx+\eta$ with $\eta\sim N(0,R)$ independent of the forecast state and SPD $R$, the likelihood is
\begin{equation}
\ell(y|x)=(2\pi)^{-m/2}|R|^{-1/2}\exp\{-\tfrac12(y-Hx)^TR^{-1}(y-Hx)\}.
\end{equation}
The evidence $Z=\int\ell p$ is positive when the prior is a nonzero nonnegative law. Equation~\eqref{eq:bayes-overview} follows from the joint density and conditional-density definition. The constant likelihood prefactor cancels in the normalized posterior, but the remaining evidence is still required.

The implemented observation choices include full state, $x$ only, and $(x,z)$. A $z$-only synthetic likelihood is used to construct the mature certification law. Partial observation can retain ambiguity; a Gaussian closure would remove multimodality by design. The density solver propagates the actual represented posterior instead.

\section{Projected analysis is a new approximation}
Let $q=\ell p_h$. The DG analysis solves $(q_h,v_h)=(q,v_h)$ for all $v_h$ using quadrature-controlled $L^2$ projection. Constants in the space imply preservation of the quadrature-defined integral before normalization. The positive product $q$ can project to a signed $q_h$, because an orthogonal projection onto a polynomial space does not preserve order. The same average repair and positivity machinery must therefore be applied after normalization or in the implementation's declared sequence.

Nodal multiplication is cheaper but differs mathematically: multiplying values at FE nodes creates an interpolant of the product, whose integral need not agree with the evidence integral. The repository retains it as an explicitly labelled comparator. Quadrature degree 14 was adopted to reduce errors for Gaussian laws and likelihood products, whose integrands are not polynomials.

\section{Analytic conditioning of Gaussian mixtures}
For component $N(m,C)$ and a linear Gaussian observation, completing the square gives
\begin{equation}
C^+=(C^{-1}+H^TR^{-1}H)^{-1},\quad
m^+=C^+(C^{-1}m+H^TR^{-1}y).
\end{equation}
The equivalent Kalman form is $m^+=m+CH^T(HCH^T+R)^{-1}(y-Hm)$ and $C^+=C-CH^T(HCH^T+R)^{-1}HC$. Mixture weights are multiplied by the component evidences $N(y;Hm,HCH^T+R)$ and renormalized. Conditioning therefore preserves a finite Gaussian-mixture form without replacing it by one Gaussian.

The mature experiment fits and mirrors sixteen components, then applies the $z=25$, variance-16 likelihood analytically. This creates one reproducible non-Gaussian posterior. It is a test law assembled from an independent spin-up ensemble, not evidence that every operational analysis posterior has the same widths, smoothness or refinement support.

\section{Saved density rather than an imagined shape}
Figure~\ref{fig:mature-density} uses the saved broader-support initial and final arrays. Integrating out one coordinate gives each joint marginal. The logarithmic colour scale makes tails visible but its $10^{-12}$ plotting floor is a display choice, not solver clipping. The figure documents the law actually used. A 2-D marginal cannot establish its complete 3-D structure, but complements the exact mass and array evidence.

\begin{figure}[p]\centering\includegraphics[width=\textwidth]{density_marginals.pdf}
\caption{Joint marginal densities from the saved broader-support MFEM run, initial projected mature law (top) and $t=0.05$ forecast (bottom). Axes are nondimensional Lorenz coordinates; colour represents $\log_{10}$ density with a display floor of $10^{-12}$. Source: \code{runs/mfem-broader-support/20261002T080644Z/mfem/*averages.bin}.}\label{fig:mature-density}
\end{figure}
\begin{figure}[ht]\centering\includegraphics[width=\textwidth]{marginal_lines.pdf}
\caption{One-dimensional marginals of the same saved initial and final fields. Coordinate axes are overlaid for comparison; each marginal separately integrates to approximately one. These plots describe a short posterior forecast, not a long-time invariant attractor.}\label{fig:marginals}
\end{figure}

\section{Learning checkpoints}
The initial-law contract has three distinct levels: the analytic mixture, its finite-element projection, and its exported voxel averages. Across different meshes the analytic law is identical while the projected polynomials can differ. Across the uniform framework equivalence test, the same already-projected positive FE law is transferred. Across half-timestep runs, matching initial voxel hashes establish the exported initial state is unchanged. These conditions answer different questions and must remain explicit.

To check understanding, derive the posterior weight update for two equally weighted Gaussian components observed only in $z$. Determine when the posterior weights remain equal and when unequal component $z$ variances break that equality. Then compare this analytic argument with the weaker statement that the Lorenz drift is symmetric. The noise and the observation determine whether that symmetry persists.

\part{Constructing the numerical method from the probability flux}
\chapter{Approximation spaces and the competing discretisations}\label{ch:spaces}
\section{What a finite-dimensional density means}
A PDE evolves a function with infinitely many possible spatial degrees of freedom. A discretisation selects a finite space and determines evolution of its coefficients. The choice must represent the resolved density structure, conserve the appropriate integral, and support the boundary and flux conventions. No finite-dimensional approximation can manufacture information about structures finer than its resolution.

A finite-volume state consists of average densities on cells. A finite-element state consists of coefficients in local or global basis functions. Continuous Galerkin enforces agreement of traces between elements; discontinuous Galerkin allows independent polynomial values on each cell and introduces interface fluxes. These are mathematical choices about the trial space and weak balance, rather than merely framework names.

The implemented DG space on affine hexahedra is
\begin{equation}\label{eq:dg-space}
V_h=\{v\in L^2(\Omega):v|_K\circ F_K\in Q_k([0,1]^3)\text{ for each }K\},
\end{equation}
where $F_K$ is the affine cell map and $Q_k$ contains tensor polynomials with degree at most $k$ in each coordinate. $Q_2$ has $(2+1)^3=27$ local DOFs. Its monomial $x^2y^2z^2$ has total degree six, so calling it ``quadratic'' must mean tensor degree, not total-degree-$2$ polynomials on a tetrahedron.

An affine diagonal Jacobian on axis-aligned boxes preserves the tensor structure and makes all cells share the same reference mass and Bernstein transformations up to scaling. This enables the reused local QP. Arbitrary warped cells would require new mapping and positivity analysis; the current claims are explicitly tied to affine cells.

\section{Why discontinuity helps and costs}
Independence of neighbouring cell polynomials permits element-local slope or QP corrections without altering continuity constraints. Constant test functions on each cell expose a local probability balance. A high-order polynomial captures variation inside a cell, while an upwind trace determines which side supplies transported probability. These are direct responses to the transport and admissibility requirements.

DG also increases unknown count. Q1 has eight local values, Q2 has 27 and Q3 has 64. On a uniform $60\times72\times72$ Q2 mesh there are $311040$ cells and $8398080$ DOFs. Matrix storage, face work, optimization calls and global output then become substantial. Equal DOFs do not imply equal cell widths or equal ability to resolve a narrow lobe: Q1 and Q2 can allocate their unknowns very differently.

\section{One-dimensional polynomial coordinates}
At nodes $0,1/2,1$, the quadratic Lagrange basis is
\begin{equation}
\ell_0=2\xi^2-3\xi+1,\quad \ell_1=4\xi(1-\xi),\quad \ell_2=2\xi^2-\xi.
\end{equation}
These functions sum to one but individual functions can be negative. Positive nodal coefficients therefore need not give a nonnegative Q2 polynomial. For example $q(\xi)=(\xi-1/4)^2-1/100$ is positive at $0,1/2,1$ and negative at $1/4$. In contrast Q1 basis functions $1-\xi$ and $\xi$ are nonnegative and partition unity, so nonnegative Q1 nodal values imply whole-interval positivity.

The exact one-dimensional average weights are $a=(1/6,2/3,1/6)$. Tensor products give $a_i a_j a_k$ in 3-D. Those weighted nodal sums provide the true average, rather than an unweighted mean of 27 Q2 coefficients. The local mass matrix is derived in Appendix~\ref{app:local-algebra}.

\section{Projection versus interpolation}
Interpolation matches selected point values. $L^2$ projection matches weighted integrals against every basis function. If $p_0$ is an analytic Gaussian or mixture, its FE projection $u_h$ solves
\begin{equation}\label{eq:l2projection}
\int_Ku_hv_h\dd x=\int_Kp_0v_h\dd x\quad\text{for all }v_h|_K\in Q_k.
\end{equation}
Taking $v_h=1$ preserves the quadrature-defined cell mass, but does not imply positivity. Gaussian integrands are nonpolynomial, so increasing integration order matters. The project separately measures raw initialization correction; it should not be counted as a dynamical positivity failure.

\evidence{\code{FokkerPlanckSolver.gaussian_projected}, \code{project_expression} and the MFEM order-14 \code{DomainLFIntegrator} implement projected initialization. Earlier nodal paths remain controlled comparators.}

\chapter{Finite volumes, exponential fitting and the reference boundary}\label{ch:fv}
\section{Deriving cell-average evolution}
For a cell $K_i$ of volume $V_i$, integration of~\eqref{eq:balance} gives $V_i\dot p_i=-\sum_{F\subset\partial K_i}|F|\widehat J_F\cdot n_i$. In one dimension a shared flux $J_{i+1/2}$ produces
\begin{equation}
\dot p_i=-\frac{J_{i+1/2}-J_{i-1/2}}{h_i}.
\end{equation}
Internal fluxes telescope in the volume-weighted sum. This conservation mechanism is independent of whether the chosen flux is accurate or positive. A single mistaken sign can conserve a wrong model, which is why consistency and comparative tests accompany mass tests.

The repository reference uses upwind drift and centred diagonal diffusion. For speed $a_{i+1/2}$ and scalar $d\ge0$,
\begin{equation}\label{eq:fv-flux}
J_{i+1/2}=a_{i+1/2}^+p_i+a_{i+1/2}^-p_{i+1}-d\frac{p_{i+1}-p_i}{h},
\end{equation}
where $a^+=\max(a,0)$ and $a^-=\min(a,0)$. Boundary face flux arrays are zero. The three-dimensional update adds directional flux differences; the reference rejects general off-diagonal $D$. It therefore supplies an independent diagonal-noise comparator but does not validate the final mixed-diffusion operator by itself.

\section{Positivity under a CFL restriction}
An explicit Euler update for~\eqref{eq:fv-flux} can be written as a nonnegative combination of nearby averages when $\Delta t$ times the total outgoing transport/diffusion rate is at most one. The off-diagonal transfer rates are nonnegative; the diagonal coefficient is $1-\Delta t\,r_i$. This gives a direct positivity proof. SSPRK(3,3) composes convex combinations of Euler steps, retaining the property under its admissible step bound.

The reference computes an allowed timestep and subdivides the interval accordingly. It checks mass and minimum density after propagation. Its diagonal covariance diagnostic adds $h_i^2/12$ to account for uniform mass within each piecewise-constant cell. Other common-grid comparators use voxel-centre moments without that addition; one must compare the same convention or account for the difference.

\section{Why exponential fitting was considered}
For $J=ap-dp'$ with constant $a,d$ on a cell interval, solving the steady local ODE gives an exponential profile. With $w=ah/d$, the exact endpoint flux can be written
\begin{equation}\label{eq:sg}
J_{i+1/2}=\frac d h\big(\mathcal B(-w)p_i-\mathcal B(w)p_{i+1}\big),\qquad
\mathcal B(w)=\frac{w}{e^w-1}.
\end{equation}
The continuous limit $\mathcal B(0)=1$ follows by Taylor expansion. This fitted flux balances advection and diffusion and approaches upwind behaviour at large positive or negative $w$. Chang--Cooper and Scharfetter--Gummel families exploit related ideas~\cite{chang}. They are relevant alternatives when equilibrium structure and directional diffusion can be used.

The Lorenz full-SPD problem contains mixed derivative terms and a rotational, compressible nonlinear drift. A directional fit applied independently to the diagonal entries would leave those mixed terms untreated. A multidimensional monotone tensor construction requires additional geometry and analysis. The method-selection report considered such alternatives, but the currently tracked finite-volume source is the simpler flux~\eqref{eq:fv-flux}. No evidence supports presenting an implemented Chang--Cooper solver as an ancestor of the final MFEM method.

\section{Context is not development history}
The distinction is worth stating because a monograph may legitimately teach a technique the project did not run. Here exponential fitting is explanatory context and a researched candidate. The actual reference contributes diagonal-law comparisons and tests. The final high-order branch emerged from same-mesh Q2 evidence and positivity diagnosis, rather than from a confirmed implementation of a directional exponential-fit branch.

\chapter{Weak balance and conservative DG advection}\label{ch:weak}
\section{Integration by parts with local traces}
Multiply $p_t+\divg(fp-D\grad p)=0$ by a test polynomial $v$ on a cell $K$. Integrating the total flux once gives
\begin{equation}\label{eq:cell-weak}
(p_t,v)_K-(fp,\grad v)_K+(D\grad p,\grad v)_K
+\langle J\cdot n,v\rangle_{\partial K}=0.
\end{equation}
The volume signs follow directly from the divergence theorem. On an interior face $F=\partial K^+\cap\partial K^-$ choose $n=n^+=-n^-$. Define scalar jump $\jump v=v^+-v^-$, vector normal jump $\jump{vn}=v^+n^++v^-n^-=\jump v\,n$, and arithmetic average $\avg q=(q^++q^-)/2$.

For continuous drift, $a_F=f\cdot n$ has a unique physical value on the face. A single advective numerical flux $\widehat f_p\cdot n$ contributes $\int_F(\widehat f_p\cdot n)\jump v$. With $v=1$ globally the jump vanishes; with $v$ equal to one on one cell and zero elsewhere, the shared flux appears as that cell's boundary exchange.

\section{Deriving upwind from characteristics}
For the one-dimensional advection equation $p_t+ap_x=0$, the characteristic arriving at a face comes from the left when $a>0$ and the right when $a<0$. The interface Riemann solution therefore supplies $p^{\rm up}=p^+$ for $a_F\ge0$ under the orientation chosen above, and $p^-$ otherwise. The project form is
\begin{equation}\label{eq:adv-form}
a_{\rm adv}(u,v)=-\sum_K\int_Kfu\cdot\grad v\dd x
+\sum_{F\in\mathcal F_h^i}\int_Fa_Fu^{\rm up}\jump v\dd s.
\end{equation}
This is precisely the Python \code{adv_volume} plus \code{adv_face}. The exterior contribution is zero total flux, so there is no independently imposed exterior upwind term in the combined formulation.

The equivalent identity
\begin{equation}
a_Fu^{\rm up}=a_F\avg u+\tfrac12|a_F|\jump u
\end{equation}
separates a central flux from dissipative jump control. Switching the normal changes $a_F$ and the scalar jump signs together, leaving the physical bilinear form invariant. A face-orientation bug often breaks this invariance while still giving apparently plausible fields.

\section{Energy and the compressible drift}
Take $u=v$ and integrate the volume term:
\begin{equation}
-\int_Kfu\cdot\grad u=-\tfrac12\int_{\partial K}(f\cdot n)u^2+\tfrac12\int_K(\divg f)u^2.
\end{equation}
Internal central-flux terms cancel the corresponding trace terms. The remainder includes $\frac12\int_F|a_F|\jump u^2$ and the compressibility contribution. For Lorenz $\divg f<0$, density $L^2$ can grow even while mass is conserved. One must not claim an unconditional $L^2$ contraction from upwinding alone. Boundary and diffusion terms determine the full energy estimate.

This distinction connects directly to probability concentration: contraction of state-space volume can raise density peaks. An increasing $L^2$ norm is not automatically a numerical instability. Instability assessment requires comparison to the correct continuous and discrete estimates, along with mass, residuals and distributional evidence.

\section{A conservation proof for the spatial form}
\begin{proposition}\label{prop:spatial-mass}
Assume the mesh integration and face assembly use shared conservative fluxes, constants belong to $V_h$, and exterior total flux is zero. Then the DG spatial form used in this project satisfies $a_h(u,1)=0$ for every $u\in V_h$ in exact arithmetic.
\end{proposition}
\begin{proof}
In~\eqref{eq:adv-form}, $\grad1=0$ and $\jump1=0$. In the diffusion form derived next, every volume or face term contains $\grad1$ or $\jump{1n}$, so it also vanishes. There is no exterior total-flux contribution. Summing therefore gives zero independently of the polynomial coefficients.
\end{proof}
This statement is about the left mass functional of the evolution matrix. It does not say a constant density is stationary under Lorenz drift. Indeed $-\divg(f\cdot1)=41/3$ in the interior. Testing mass preservation and testing operator action on a constant are different operations.

\chapter{Tensor diffusion and the SIPG construction}\label{ch:sipg}
\section{Why volume diffusion alone is incomplete}
For discontinuous $u_h$, a cellwise gradient omits jumps across faces. Summing only $\int_K\grad u^TD\grad v$ would not recover the correct weak flux coupling. Consistency requires a numerical normal diffusive flux. Symmetry and control of jumps supply the remaining SIPG terms.

The project diffusion form is
\begin{align}\label{eq:sipg}
a_D(u,v)={}&\sum_K\int_K(D\grad u)\cdot\grad v\dd x\\
&-\sum_F\int_F\avg{D\grad u}\cdot\jump{vn}\dd s\nonumber\\
&-\sum_F\int_F\avg{D\grad v}\cdot\jump{un}\dd s\nonumber\\
&+\sum_F\int_F\eta_F\jump{un}\cdot\jump{vn}\dd s.\nonumber
\end{align}
The first face term is obtained from the integration-by-parts diffusion trace. The second is its adjoint counterpart, restoring symmetry. The final penalty controls interelement jumps. For a globally smooth $u$, $\jump{un}=0$ and the last two terms involving that jump vanish; the remaining form reproduces integration by parts with its physical diffusive flux.

The code uses
\begin{equation}\label{eq:penalty}
\eta_F=\gamma k^2\frac{\avg{n^TDn}}{\avg{h_K}},\qquad\gamma=16,
\end{equation}
with cell diameter $h_K$ and tensor degree $k=2$, giving the factor 64. This is the actual penalty convention; replacing diameter by normal cell width would change the method. MFEM's default face scaling needed a custom integrator to match it physically.

\section{Coercivity under a sufficiently large penalty}
For symmetric $D\succ0$, define a broken energy seminorm
\begin{equation}
\|v\|_{DG,D}^2=\sum_K\|D^{1/2}\grad v\|_{L^2(K)}^2+\sum_F\|\eta_F^{1/2}\jump v\|_{L^2(F)}^2.
\end{equation}
An inverse trace estimate bounds normal-gradient traces by a mesh- and degree-dependent multiple of cell gradient energy. Apply Cauchy--Schwarz and Young's inequality $2ab\le\epsilon a^2+\epsilon^{-1}b^2$ to the two consistency terms in $a_D(v,v)$. If the penalty dominates the resulting trace constant, then
\begin{equation}
a_D(v,v)\ge c\|v\|_{DG,D}^2
\end{equation}
modulo the constant nullspace for no-flux diffusion. Shape regularity, tensor ellipticity and correctly scaled penalty are required. Arnold and colleagues provide a unifying framework for DG elliptic analysis~\cite{arnold}. That framework explains the construction; it does not prove that the fixed numerical factor 16 is sufficient on every hypothetical distorted or extremely graded mesh.

For affine rectangular meshes, the code's scaling is controlled empirically through equivalence, mass and diffusion probes. A rigorous global coercivity certificate for the complete NC hierarchy has not been recorded. The text must distinguish that open proof from successful numerical propagation.

\section{What anisotropy changes}
The energy uses $\grad v^TD\grad v$, so the diffusion directions are the eigenvectors of $D$, not necessarily the coordinate axes. The face quantity $n^TDn$ captures normal diffusion strength. The cross terms affect both volume gradients and numerical face fluxes. The matrix-valued MFEM integrator accepts this tensor, but compatibility of API types is weaker evidence than matched operator actions.

Very small $\lambda_{\min}(D)$ can create directionally weak smoothing and high directional P\'eclet numbers. Semidefinite $D$ requires a different coercivity and regularity analysis. The chosen test tensor is moderately anisotropic and uniformly SPD; general SPD parameterization for a pilot must not silently inherit its numerical guarantees for nearly singular matrices.

\section{Exterior reflecting treatment}
Equation~\eqref{eq:cell-weak} has a \emph{combined} boundary term. After replacing interior traces by upwind and SIPG fluxes, the exterior combined flux is prescribed zero. The formulation thus has no independent Dirichlet penalty, prescribed diffusion Neumann datum or exterior advective inflow law. A separate pure-diffusion operator diagnostic has natural homogeneous Neumann behaviour, but that isolated diagnostic is not the full reflecting convection--diffusion boundary model.

\evidence{Python lines defining \code{sipg} in \code{lorenz_fpe/core.py}; MFEM \code{DolfinxPenaltyDGDiffusionIntegrator}; manufactured boundary convergence in \code{boundary_flux_report.json}.}

\chapter{From weak forms to matrices and Crank--Nicolson}\label{ch:cn}
\section{The sign convention that must survive a port}
Let $p_h=\sum_ju_j\phi_j$. With $M_{ij}=(\phi_j,\phi_i)$ and $S_{ij}=a_h(\phi_j,\phi_i)$, the semidiscrete equation is
\begin{equation}\label{eq:semi}
M\dot u+Su=0,\qquad M\dot u=K_hu,\quad K_h=-S.
\end{equation}
$M$ is symmetric positive definite for an independent basis and positive cell measures. It is block diagonal by cell for DG, though implementations may assemble distributed sparse matrices. $S$ is nonsymmetric because of upwind drift. Its diffusion part is symmetric under matched SIPG conventions.

Python assembles $M+\theta\Delta tS$; the MFEM evolution matrix is $K_h$ and its CN left matrix is $M-\frac12\Delta tK_h$. Both describe the same equation. Comparing source-level signs without identifying which matrix represents spatial balance can produce a false mismatch or mask a true one. The early MFEM probe did discover a real sign defect, caught before enabling AMR.

\section{Derivation of the theta update}
Integrating~\eqref{eq:semi} over $[t_n,t_{n+1}]$ and approximating the right side with endpoint weights gives
\begin{equation}\label{eq:theta}
(M-\theta\Delta tK_h)u^{n+1}=(M+(1-\theta)\Delta tK_h)u^n.
\end{equation}
$\theta=1$ is backward Euler. $\theta=1/2$ is Crank--Nicolson. Taylor expansion for sufficiently smooth $u$ gives local defects of order $\Delta t^2$ for backward Euler and $\Delta t^3$ for CN, corresponding to first and second global order on a fixed finite-dimensional system when stability and consistent solves hold.

This is a formal time-integrator statement. The executed method replaces the raw result $u^{n+1,*}$ by a corrected density $\mathcal P(u^{n+1,*})$. Its convergence depends on the size, consistency and accumulation of that correction. The CN derivation cannot by itself prove second order for the corrected trajectory.

\section{Stability and a counterexample to positivity}
For $\dot y=\lambda y$, CN has amplification $R(z)=(1+z/2)/(1-z/2)$ with $z=\lambda\Delta t$. If $\operatorname{Re}z\le0$, then $|R(z)|\le1$, showing scalar A-stability. As $z\to-\infty$, $R(z)\to-1$, so stiff modes are not strongly damped. Backward Euler has $R(z)=1/(1-z)$ and tends to zero in that limit.

The distinction matters for positivity. Even $\dot y=-ay$ with $a>0$ and $y_0>0$ produces a negative CN result if $a\Delta t>2$. A stable high-order raw update can consequently contain negative polynomial structure. Implicitness controls one restriction, while admissibility remains a separate algebraic requirement.

\section{Mass with exact and inexact linear solves}
Let $e$ represent the constant test function. Then the discrete integral is $m^Tu$ with $m^T=e^TM$, while $e^TK_h=0$ by Proposition~\ref{prop:spatial-mass}. Multiplying~\eqref{eq:theta} by $e^T$ proves
\begin{equation}\label{eq:cn-mass}
m^Tu^{n+1}=m^Tu^n.
\end{equation}
If the solve residual is $r=b-A_{CN}u^{n+1}$, the mass defect satisfies
\begin{equation}
m^T(u^{n+1}-u^n)=-e^Tr
\end{equation}
under that residual sign convention. Thus tolerances and cancellation affect measured conservation. An absolute accumulated action such as $-5.6\times10^{-10}$ is uninterpretable without scale; a normalized action residual of $1.93\times10^{-17}$ supports machine-precision structural conservation in the production probe.

\section{Nonsymmetric solves and reused operators}
GMRES builds a Krylov subspace and minimizes a residual norm for nonsymmetric systems. A preconditioner changes the conditioning of the iteration, not the intended solution. The Python solver uses ILU in serial and block Jacobi under MPI, with independent true residual evaluation. MFEM uses Hypre GMRES with ILU(1), restart dimension 100, up to 1000 iterations, relative tolerance $10^{-14}$ and absolute tolerance $10^{-17}$ in the mature branch.

The mature MFEM loop does not write a full per-step true-residual history in its summary. Tight requested tolerances and equivalence are strong practical evidence, but they do not substitute for a separately recorded residual gate on every trajectory. This reporting limitation is stated in the final specification.

Constant coefficients, a fixed mesh and fixed timestep allow reuse of CN matrices and preconditioners. Changing mesh mid-forecast would require rebuilding or transferring operators and states. The project first chooses a static AMR mesh, removing this additional source of error and cost from its equivalence study.

\chapter{Whole-cell positivity and Bernstein certificates}\label{ch:bernstein}
\section{Several notions that must not be conflated}
Nonnegative integral mass, nonnegative cell averages, nonnegative nodal values, nonnegative sampled values and a nonnegative polynomial everywhere are distinct properties. Positive averages constrain one linear functional per cell. Q2 has 26 additional mean-free directions that can produce internal oscillations. A certificate must control those directions on the entire cell or declare its coverage incomplete.

The Bernstein basis on $[0,1]$ is $b_i^k(\xi)=\binom ki\xi^i(1-\xi)^{k-i}$. Each basis function is nonnegative and $\sum_ib_i^k=1$. Tensor products retain both properties. For coefficients $\beta_{ijk}$,
\begin{equation}\label{eq:bernstein-enclosure}
\min\beta_{ijk}\le p(\xi,\eta,\zeta)\le\max\beta_{ijk}\quad\text{on }[0,1]^3.
\end{equation}
This enclosure follows because the polynomial value is a convex combination of coefficients. Therefore coefficient nonnegativity is a sufficient whole-cell certificate.

\section{A sufficient condition with an incomplete converse}
A nonnegative polynomial can have negative Bernstein coefficients. For $p(x)=(x-c)^2$ on an interval $[a,b]$ containing $c$, the middle quadratic Bernstein coefficient is $(a-c)(b-c)<0$. The endpoint coefficients are squares. If $c$ is irrational, no finite dyadic subdivision puts its zero exactly on a subdivision boundary. The interval containing $c$ retains a negative middle coefficient at every finite depth.

This gives a precise counterexample to completeness of recursive positivity certification. Failure to certify may mean genuine negativity or an enclosure that straddles zero. The implementation separates three outcomes: certified nonnegative, witnessed negative by point evaluation beyond scale-aware tolerance, and unresolved at the depth limit. Both latter outcomes require a conservative correction in the propagated method.

\section{Subdivision and the actual constraint set}
de Casteljau subdivision expresses the same polynomial in Bernstein form on smaller cells. New coefficients are convex combinations of old ones, so enclosure bounds tighten. The project's fixed Q2 constraint matrix uses two subintervals per axis, giving eight control subcells with 27 coefficients each, or 216 rows before duplicate removal. The optimizer keeps a deduplicated matrix, while acceptance is checked against the full original rows.

Adaptive depth-four subdivision is used to skip a polynomial that is already certifiably nonnegative. It does not alter the convex feasible set while the QP is being solved. The mature diagnostic compares depth four and eight on the same saved raw polynomials. This prevents changes in dynamics from being mistaken for changes in certification.

\section{Floating-point qualification}
Exact-arithmetic Bernstein coefficients give an exact sufficient condition. Computed transformations, mapping and optimizer output contain roundoff. The code uses scale-aware witness tolerances and a normalized positivity margin $10^{-12}$ for QP constraints, then rechecks rows and averages. This is practical numerical certification, not outward-rounded interval arithmetic. The book therefore reports ``whole-cell Bernstein-certified at the implementation tolerance'' rather than asserting a computer-assisted exact positivity proof.

Cells with nonpositive repaired average are set to the zero polynomial. For a nonnegative continuous polynomial, zero integral implies the polynomial is zero everywhere: if it were positive at one point, continuity would give a neighbourhood of positive integral. That algebraic fact explains the zero-average branch. Tiny negative average values admitted by numerical tolerances must still be controlled by global mass checks.

\section{Why deeper certificates were rejected}
The mature experiment found that depth eight rescued at most 0.0788 percent of depth-four noncertified cells and reduced correction by at most 2.89 percent. Direct negative witnesses dominated. A tighter enclosure can reveal positivity hidden by a sufficient condition; it cannot change a genuinely negative polynomial. The diagnostic therefore narrowed the relevant hypothesis from certificate conservatism towards actual approximation or update error.

\evidence{\code{mature_positivity_certificate_diagnostic_report.json}, \code{mature_positivity_diagnostic.py}, \code{PositivityLimiter._classification_summary} and MFEM \code{AdaptiveCertificate}.}

\chapter{Global average repair and local minimum-change QP}\label{ch:qp}
Optimisation-based Fokker--Planck limiters provide a neighbouring two-stage architecture~\cite{liuhu}. The published semi-implicit temporal/flux construction differs from this fully CN/upwind/SIPG implementation. Its theorem is therefore not invoked as a certificate for the present three-dimensional forecast.
\section{The nonnegative equal-mass projection}
For raw cell averages $w_K$ and volumes $V_K>0$, Stage 1 solves
\begin{equation}\label{eq:stage1}
\min_{x_K\ge0}\frac12\sum_KV_K(x_K-w_K)^2,
\quad\sum_KV_Kx_K=\sum_KV_Kw_K=M_0.
\end{equation}
It is feasible if and only if $M_0\ge0$. When feasible, strict convexity gives a unique minimizer. A Lagrange multiplier $\lambda$ for mass and multipliers $\nu_K\ge0$ for the lower bounds yield $V_K(x_K-w_K+\lambda)-\nu_K=0$ and $\nu_Kx_K=0$. Thus
\begin{equation}\label{eq:waterfill}
x_K=\max(0,w_K-\lambda).
\end{equation}
The scalar function $F(\lambda)=\sum_KV_K\max(0,w_K-\lambda)$ is continuous and nonincreasing. For $M_0>0$ it is strictly decreasing near the required positive active set, yielding a unique multiplier there. Bisection or a sorted active-set calculation can recover it. $M_0=0$ yields the unique vector zero, while the multiplier need not be unique.

Weighted mass is essential on NC meshes. If one treats a small refined cell and a large background cell as equal-volume entries, conserving their unweighted sum conserves the wrong quantity. The MFEM average repair uses actual transformation volumes and MPI reductions; the Python baseline gathers averages/volumes to rank zero and scatters the projected result.

The cell polynomial is shifted by the constant $x_K-w_K$. Since the basis represents one exactly, this changes its average without altering its mean-free shape. It makes the subsequent local positivity problem feasible.

\section{Scalar contraction as a feasible comparator}
If a cell has corrected average $\bar p\ge0$ and sufficient lower enclosure $m_K$, contraction gives
\begin{equation}\label{eq:scaling}
p^{\rm scale}=\bar p+\theta(p-\bar p),\quad
\theta=\min\left(1,\frac{\bar p}{\bar p-m_K}\right)
\end{equation}
when $m_K<0$, with the zero-average case treated separately. Its average is unchanged because the mean-free component integrates to zero. If all corrected Bernstein coefficients are nonnegative, it certifies the polynomial. The operation scales every higher mode by the same factor, which can erase accurate shape information away from the offending direction.

The project uses scalar scaling as a safe feasible candidate, diagnostic comparator and objective bound for the local QP. A minimum-change convex projection cannot have larger mass-matrix objective than that feasible candidate when solved exactly. This is an objective-norm statement; it does not automatically imply a smaller $L^1$ correction or better covariance. Those comparisons require measurements.

\section{The local constrained problem}
Let $c_0\in\R^{27}$ be the raw Q2 nodal vector after Stage 1. Let $H$ be the SPD reference mass matrix, $a^Tc$ its cell average and $Cc$ its fixed-subcell Bernstein coefficients. The local problem is
\begin{equation}\label{eq:local-qp}
\min_c\tfrac12(c-c_0)^TH(c-c_0),\qquad Cc\ge0,\quad a^Tc=a^Tc_0.
\end{equation}
$H$ makes the objective the squared $L^2$ change on the reference cell. On affine cells the physical mass matrix is $V_KH$, and multiplying the objective by a positive $V_K$ does not change its minimizer. This is the mathematical reason the same reference QP can be reused throughout a static locally refined mesh.

\begin{proposition}[Feasibility and uniqueness]\label{prop:qp}
If the prescribed average $\bar p$ is nonnegative, constants are represented by $\bm1$, $a^T\bm1=1$ and $C\bm1=\bm1$, then~\eqref{eq:local-qp} has a unique solution. If $\bar p<0$, no nonnegative polynomial can meet that mass constraint.
\end{proposition}
\begin{proof}
For $\bar p\ge0$, $c=\bar p\bm1$ is feasible. The feasible set is closed and convex, and the SPD quadratic objective is coercive and strictly convex. Existence follows by restricting a minimizing sequence to a bounded sublevel set and compactness in finite dimension; strict convexity gives uniqueness. Negative average contradicts nonnegativity of the polynomial integral. For $\bar p=0$, the only nonnegative polynomial is zero by continuity and the integral argument.
\end{proof}

\section{Normalization and equality elimination}
For $\bar p>0$, set $z_0=c_0/\bar p$ so $a^Tz_0=1$. The positive margin is applied in this normalized space, avoiding ill-conditioned absolute constraints in low-density tails. Let $N\in\R^{27\times26}$ have orthonormal columns spanning $\ker a^T$. Write $z=z_0+Ny$. Then mass is automatic and the QP reduces to
\begin{equation}\label{eq:reduced-qp}
\min_y\tfrac12y^T(N^THN)y,\qquad CNy\ge\varepsilon\bm1-Cz_0,
\end{equation}
where $\varepsilon=10^{-12}$. The reduced Hessian and constraint matrix are fixed; only the lower bound changes from cell to cell. Factorization caching and warm starts make OSQP a suitable implementation mechanism~\cite{osqp}. That external solver does not itself prove polynomial positivity; the project checks the original constraints afterwards.

Python computes a nullspace with SciPy. MFEM constructs a Householder basis. The reduced coordinates may differ while representing the same constrained physical polynomial. This is another reason raw coefficient or internal-QP variable comparison across frameworks is weaker than a physical-field comparison.

\section{Acceptance and numerical guards}
Python retains SLSQP as an oracle. MFEM accepts OSQP solved or solved-inaccurate status, then reconstructs the candidate, removes average drift by a constant shift, and if needed performs a tiny contraction toward the unit-average constant. It checks all 216 rows and marks certification only when the final minimum is finite and nonnegative. Those guard operations mean the floating-point returned state may differ slightly from the exact QP minimizer.

The reported zero optimizer failures and zero fallbacks on completed MFEM runs should be read with the source. Optimizer errors use aborting verification checks; successful summaries write zero failure counts. A run that aborts before its summary cannot be counted as a certified zero-failure forecast. Aborting failed solves is a safeguard, while richer attempted-solve histories would improve forensic reporting.

\section{What repeated correction can reveal}
Define $C_n=\|p_h^{n+1}-p_h^{n+1,*}\|_{L^1}/|M(p_h^{n+1,*})|$. Small $C_n$ supports a low-intervention interpretation. Large values can come from inaccurate initialization, temporal oscillations, an overly restrictive feasible set, poor spatial resolution or genuine physical concentration represented in an inadequate space. Correction by itself does not identify the cause.

The project separates these possibilities by frozen-polynomial certificate tests, same-positive-initial-state AFC comparisons, one-step DOF experiments and refinement. The convergence of a norm-constrained projection on one state cannot replace convergence of the full propagated density. Repeated local changes may accumulate or cancel; both per-step and full-density evidence are required.

\chapter{Alternative positivity architectures and what they tested}\label{ch:afc}
\section{A positive low-order update}
An invariant-domain architecture starts from a low-order state that is already admissible, then adds limited high-order corrections. For cell averages on a graph, a positive update has nonnegative transfer rates and a timestep bound on outgoing probability. The mature AFC comparator decomposes the chosen full tensor into axis and pair-diagonal graph directions. It requires nonnegative remaining diagonal rates after accounting for mixed entries; arbitrary SPD matrices need not admit that particular sparse decomposition on a chosen mesh.

The high-order target remains the frozen Q2 CN result. The graph step provides admissible averages; conservative antidiffusive changes move those averages towards the target. This changes the update architecture at the average level while retaining the polynomial QP afterwards. It is a targeted diagnostic of whether average negativity is the dominant mature bottleneck.

\section{The implemented complete-graph correction}
Let $l_i$ be low-order averages and $h_i$ a high-order target rescaled to the same sum. Put $\delta_i=h_i-l_i$. Negative corrections are limited by $\delta_i\ge-l_i$, and positive corrections are multiplied by a common factor so that accepted gain equals accepted loss. On the uniform comparator mesh, equal cell volumes make sum conservation equivalent to mass conservation. This is an experiment-specific complete-graph decomposition, not the same edge-local scheme as every published AFC method.

The distinction limits what a negative result says. The general literature provides conservative limiter constructions for scalar transport~\cite{kuzmin}; it does not supply a ready-made theorem that this full three-dimensional Lorenz/Q2/SIPG/CN variant will reduce correction. Carrillo--Liu--Yu explicitly separate average positivity from subsequent polynomial limiting for their Fokker--Planck class~\cite{carrillo}. The project uses that distinction as a diagnostic architecture and does not import their different energy-dissipation theorem.

\section{Why almost complete antidiffusion acceptance matters}
The mature experiment accepted more than 99.9998 percent of the high-order average correction while still requiring a local-QP change of order $10^{-3}$. This supports the conclusion that cell averages were already essentially admissible. Within-cell high-order modes remained negative. Improving average-level positivity could therefore leave the principal correction nearly unchanged.

The result is specific to that comparator and law. It does not prove every AFC or subcell method ineffective. It does show that the tested average-level intervention addressed a component with little remaining error, while the expensive polynomial repair persisted. That is enough to stop spending substantial compute on another identical average-level hierarchy.

\section{The Bernstein-DOF diagnostic}
The next prototype converts each whole-cell Q2 polynomial to 27 Bernstein coefficients. All Bernstein basis functions have integral $1/27$ on the cube, so coefficient-sum conservation enforces its cell average. A Euclidean simplex projection seeks $\beta_i\ge\epsilon$ with $\sum_i\beta_i=27\bar p$. It then transforms back to nodal coordinates.

This minimizes a coefficient-space Euclidean norm, whereas the retained QP minimizes a physical mass-matrix norm and uses more permissive subcell constraints. The objectives and feasible sets differ. Whole-cell Bernstein positivity can be stricter than positivity after subdivision. The prototype is thus a diagnostic convex coefficient repair, not an implemented sparse subcell flux-limited DG propagation operator.

The recorded comparison failed the proposed factor-of-two reduction against the local QP. A uniform-$60$ one-step check reduced QP correction substantially compared with $45$, supporting under-resolution as the stronger hypothesis. The failure catalogue must describe the prototype precisely: it does not license a general conclusion that every DOF-level invariant-domain method is inferior.

\section{Retaining the least disruptive known admissibility method}
By this point the local QP had remained the best tested repair while mesh refinement reduced both raw negativity and intervention. The next change therefore targeted resolution, with positivity mathematics frozen. This preserved a clear comparator and avoided introducing polynomial order, different fluxes and different boundary treatment simultaneously.

\chapter{Convergence, export and the meaning of a density comparison}\label{ch:convergence}
\section{Consistency, stability and convergence}
Consistency compares the discrete equation with the continuous equation when an exact sufficiently smooth solution is inserted. Stability controls amplification of perturbations and errors in the chosen norm. Convergence compares a sequence of discrete solutions with the continuum solution as resolution is refined. These ideas are linked under theorem-specific hypotheses; a bare invocation of consistency plus stability is insufficient for the nonlinear corrected scheme without identifying spaces, norms and uniform constants.

If a smooth solution has approximation error bounded by $Ch^{k+1}$ in $L^2$, the constant depends on derivatives that can become large for narrow mature structures. A nominal Q2 approximation estimate therefore does not say a particular practical mesh is accurate. The measured initialization corrections $0.380\to0.158\to0.0313$ show how far that law initially lay from the positive representable set on early meshes.

\section{Observed order from adjacent differences}
For constant ratio $r>1$, suppose $p_h=p+ch^s+o(h^s)$ in a suitable norm with a common nonzero leading error direction. Then adjacent differences approximately satisfy $D_{h,h/r}/D_{h/r,h/r^2}\approx r^s$, giving
\begin{equation}\label{eq:observed-order}
s_{\rm obs}=\frac{\log(D_1/D_2)}{\log r}.
\end{equation}
The asymptotic expansion and matched law/time/domain assumptions are necessary to interpret this as an order estimate. Without them the same formula is simply a ratio diagnostic. Different nonuniform support/depth designs lack a single scalar refinement ratio; a rate assigned to them is heuristic.

The startup four-level temporal differences $5.238\times10^{-4}$, $5.418\times10^{-4}$ and $1.620\times10^{-4}$ give orders approximately $-0.049$ and $1.741$. The plateau precedes a fine-level decrease. The final value supports entry into a better-resolved temporal regime, but does not prove uniformly second-order corrected CN convergence.

Identity-noise spatial differences $0.141394$ and $0.0355845$ at ratio $1.5$ give observed common-grid rate $3.4026$. Full-SPD differences $0.139500$ and $0.0348401$ give $3.4215$. Those rates exceed three but are measured on finite grids and projected initial states; they remain observed common-grid rates rather than formal-order proofs.

\begin{figure}[ht]\centering\includegraphics[width=\textwidth]{convergence_evidence.pdf}
\caption{Recorded startup temporal differences (left) and mature uniform-mesh mean corrections (right). The temporal horizontal coordinate labels the coarser member of each pair relative to the first timestep. Mature correction values are rounded from reports; the dashed line is the original $10^{-3}$ gate. The two panels refer to different initial laws and are not a shared convergence hierarchy.}\label{fig:convergence}
\end{figure}

\section{Conservative voxel averaging}
For export voxel $V$, define $\bar p_V=|V|^{-1}\int_Vp_h$. Splitting each voxel into intersections with FE cells permits exact integration of a Q2 polynomial by sufficient tensor Gauss quadrature on affine subboxes. The MFEM exporter uses three-point Gauss nodes and weights per coordinate, exact through degree five, hence exact for tensor Q2 on each intersection.

The global exported integral is $\sum_V|V|\bar p_V=\int p_h$ in exact arithmetic if intersections partition the cells and no region is lost or duplicated. The common $180\times216\times216$ grid lets different meshes be compared without point interpolation. Its voxel volume is $384000/8398080$.

\begin{proposition}[Averaging is an $L^1$ contraction]\label{prop:export}
Let $P_V$ map integrable functions to piecewise-constant voxel averages on a fixed partition. Then $\|P_Vp-P_Vq\|_{L^1}\le\|p-q\|_{L^1}$.
\end{proposition}
\begin{proof}
For each voxel, $|V|\,|\bar p_V-\bar q_V|=|\int_V(p-q)|\le\int_V|p-q|$. Sum over the disjoint partition.
\end{proof}
Thus the reported voxel difference is a lower bound on the true FE-polynomial $L^1$ difference, not an upper bound. Within-voxel errors can cancel. This is the precise reason subvoxel accuracy remains separate even after a small exported difference.

\section{Norms and statistical quantities}
For normalized densities, total variation is $\frac12\int|p-q|$ and equals the maximum probability discrepancy over measurable events. For bounded $g$, $|\E_pg-\E_qg|\le\|g\|_\infty\|p-q\|_1$. Moment errors are therefore controlled by full-density errors on a bounded box, while the converse fails.

Relative covariance error is a Frobenius discrepancy divided by the reference Frobenius norm. Off-diagonal correlations divide by standard deviations and can be sensitive near zero variance. Lobe probability is an event integral, here often $x>0$ versus $x<0$. Joint-TV comparisons against Monte Carlo use fixed histogram resolution and declared smoothing; they are comparisons of those processed representations rather than exact total variation of the continuum laws.

\section{Initial differences and reference limitations}
For two meshes, decompose a terminal difference conceptually into initialization, propagation and export effects. Independently projecting the same mixture does not make initial FE polynomials identical. Reporting initial voxel $L^1$ separately helps interpret the final result. Across frameworks, transfer of one projected state avoids that confound entirely.

Increasing reference particle count strengthens moments but not automatically fine density estimates. Histogram and kernel bandwidth choices introduce bias. Similarly, reproducing a uniform-$60$ reference to $0.001688$ establishes agreement with that reference, not its accuracy. The later fine graded change of $0.108680$ demonstrated why the reference question matters.

\part{The recorded research programme}
\chapter{The baseline audit: initial laws, analysis and independent references}
\section{The first hidden error was already present at time zero}
An early FE comparison used a normalized nodal Q1 interpolant while Monte Carlo sampled the whole-space analytic Gaussian. Its endpoint discrepancy included initial representation error. The audit established this confound directly, rather than attributing the entire discrepancy to PDE propagation.

On the $30\times36\times36$ Q1 mesh, the nodal initial-law density discrepancy was about 0.141348 in $L^1$, compared with about 0.0555363 for unlimited $L^2$ projection and 0.0580728 after limiting that projection. The corresponding normalized covariance discrepancies were about 0.126186, $5.10\times10^{-9}$ and 0.006835. The projected field's minimum was about $-4.91\times10^{-5}$. Its excellent moment agreement therefore coexisted with an inadmissible polynomial and appreciable density error.

The small covariance discrepancy does not establish a small full-density error. Projection preserves integrals against test functions in its space; second moments outside that space need separate analysis. The listed values are measured quantities, not a general moment-exactness theorem for Q1.

The quadrature study compared initial-projection coefficient vectors with an order-16 comparator. Relative discrepancies fell from about 0.01065 at order 4 to $5.97\times10^{-4}$ at order 6, $2.27\times10^{-5}$ at order 8, $6.14\times10^{-7}$ at order 10, $1.28\times10^{-8}$ at order 12 and $2.66\times10^{-10}$ at order 14. This justified order 14 for those laws. The order-16 field is itself a numerical comparator, and a narrower later mixture can require another quadrature check.

\evidence{The \code{initial_representation} section of \code{method_selection_report.json}, the initialization study scripts and associated comparison records provide the numbers and conventions.}

\section{The sampler must represent the law actually propagated}
A nodal or projected FE field defines its own normalized law when admissible. If a native sampler draws from that FE law, it tests propagation without silently substituting a different analytic law. Alternatively, both methods may start from one analytic law, provided FE representation error is explicitly measured. These are distinct controlled experiments.

The same-native-law Q1 forecast used $\Delta t=.000625$, horizon .05, 200,000 particles and MC timestep .000125. Covariance discrepancy about 0.0855765 was roughly 13.62 times the reference bootstrap p95 about 0.00628239. Halving the MC timestep changed normalized covariance by about 0.000110865. This isolated a substantial deterministic-method discrepancy which could not plausibly be explained by the measured MC timestep change alone. It did not provide an exact covariance-error decomposition.

Native sampling of a polynomial requires more than interpreting nodal coefficients as probabilities. The density includes basis shape and volume. A sampler must integrate or bound the nonnegative polynomial within cells, draw cell mass correctly and use an accepted local law. The corresponding source-level sampling contracts are part of the validation interface; an analytic Gaussian sampler serves a different input contract.

\section{Analysis projection replaced a nodal product}
The analytic Bayesian posterior is proportional to $\ell p$. Evaluating that product at FE nodes and interpolating it can alias a concentrated likelihood. The replacement solves the weak projection of $\ell p$, normalizes by FE evidence and applies conservative admissibility correction.

In the analytic Gaussian test on the fine Q1 mesh, nodal analysis gave mean-error norm about 0.07468, $L^1$ about 0.19361 and covariance discrepancy about 0.14494. Projected analysis improved these to about 0.002055, 0.11554 and 0.03706. Its limiter correction was still about 0.05688. The change was retained because it improved the controlled posterior representation; the remaining intervention precluded using the test as production posterior certification.

Analytic mixture conditioning offers a useful independent oracle: each Gaussian component updates by a linear-Gaussian formula, and its weight changes by its observation evidence. This preserves multimodality rather than fitting a single Gaussian. The mature mixture later used exactly this principle to separate a continuous posterior law from a mesh-specific representation.

\section{The coarse repeated-analysis studies}
The earlier \code{mature_da_report.json} evolved 20 independent trajectories for 20 cycles on $8\times10\times10$ Q1 cells, with $\Delta t=.005$, observation interval .05, observed coordinates $(x,z)$, observation variance 9 and seed 71023. The 800 cells provided only 6,400 DG unknowns. The reported corner drift speed approached 1,975 and the corresponding coarse P\'eclet estimate was approximately 15,801, indicating a strongly transport-dominated discretisation by that diagnostic.

The report's explicit classification was exploratory coarse distribution study, not ground truth. Cycle-mean relative limiter intervention ranged from about 0.06781 to 0.11344. A feature-stationarity criterion did not establish stable burn-in in the available horizon; the suggested cycle 15 was provisional.

Regeneration with projected analysis, documented in \code{mature_da_projected_report.json}, retained the trajectory/cycle counts, mesh, observation model and seed. Cycle-mean intervention ranged from about 0.06577 to 0.10214. The specified feature criterion selected burn-in 12, using the final-window comparison. This tests stability of those selected feature means, not stationarity of the entire posterior law, convergence to an invariant density or accuracy of the coarse FE trajectory.

Reported wall times were approximately 97.2 and 102.5 seconds. Their low cost reflects the coarse exploratory configuration. Those outputs remain valuable for testing data schemas, trajectory splits and analysis workflows. They cannot be relabelled as the later high-resolution mature NC-hex pilot merely because both involve data assimilation.

\section{Independent finite volumes were informative but unresolved}
The diagonal-noise independent reference refined cell-centred upwind transport and centred diffusion under SSPRK(3,3). At refinement factors 1, 2 and 4, normalized covariance discrepancy versus the same MC reference was about 0.58833, 0.27181 and 0.13084. The finest grid was $120\times144\times144$ and took about 66.0 seconds for its recorded run. Q1 versus finest FV density difference was about 0.15390.

The trend supports reference improvement, yet the remaining discrepancy is too large to call the finest FV field truth. Two numerical methods approaching each other can share error, and two methods disagreeing substantially do not identify which is correct without further evidence. The FV hierarchy remained an independent diagonal comparator with a first-order spatial transport limitation.

\section{Monte Carlo density reconstruction was not selected as ground truth}
The 200,000-particle experiment had accurate and cheap moment diagnostics but sparse occupancy on a $120\times144\times144$ output grid. Raw-histogram differences were about 0.19824 versus FV and 0.24284 versus Q1. Five batch histogram TVs from the full histogram were around 0.119--0.122, showing appreciable reconstruction variability.

Kernel smoothing reduced one comparison at selected bandwidths but introduced strong bandwidth dependence: $L^1$ versus FV was about 0.09896, 0.08899, 0.17392 and 0.30175 for bandwidths 0.75, 1.0, 1.5 and 2.0 voxels. A tail region with fewer than one expected particle per voxel carried about 0.01378 probability in the FV representation but about 0.006205 in the MC histogram. These quantities are comparisons of estimators, with FV error also present; they are evidence against treating a sparse raw histogram as a stable full-density target.

\section{Researched alternatives with no executed solver branch}
The method-selection record considered a boundary-consistent characteristic remap, translated/scaled Hermite-Galerkin references and directional exponentially fitted families. The characteristic challenger explicitly remained unimplemented because its available theory did not establish the intended reflecting Lorenz boundary and whole-cell positivity in three dimensions. A Hermite expansion could represent a whole-space smooth density economically, but spectral coefficients do not automatically define a nonnegative density; mode and domain/reference checks would still be needed.

These proposals belong in the portfolio and literature context. They are not archived successful ancestors of the final method. Likewise, the early finite-volume source uses the upwind formula derived in Chapter~\ref{ch:fv}, with no confirmed Chang--Cooper implementation. Preserving this distinction prevents a plausible textbook narrative from replacing the actual record.

\chapter{How the evidence changed the numerical question}
\section{A development history needs controlled comparisons}
A software history and a scientific history answer different questions. A commit can show when an implementation changed. It cannot, by itself, show that the new implementation solved the intended equation more accurately. A run directory can contain a completed process while its scientific decision rejects the result. The development below therefore follows a chain of hypotheses, controlled comparators, measured quantities and decisions. Appendix~\ref{app:ledger} records the source inventory; the following chronological appendix preserves all 69 available commit messages rather than inventing a history before the first recorded baseline.

At each stage, write the numerical output as
\begin{equation}
 p_{h,\Delta t,\Omega,\mathcal L,\mathcal I},
\end{equation}
where $\mathcal L$ denotes the positivity treatment and $\mathcal I$ the initial representation. A difference between two such fields can be attributed to one factor only when the other factors are fixed, or when their influence is independently bounded. This is a practical experimental-design principle rather than a claim that nonlinear numerical errors add independently.

For example, the inequality
\begin{align}
 \|p_{h_1,\Delta t_1,\mathcal L_1}-p_{h_2,\Delta t_2,\mathcal L_2}\|_1
 \le{}&\|p_{h_1,\Delta t_1,\mathcal L_1}-p_{h_1,\Delta t_2,\mathcal L_1}\|_1\\
 &+\|p_{h_1,\Delta t_2,\mathcal L_1}-p_{h_2,\Delta t_2,\mathcal L_1}\|_1
 +\|p_{h_2,\Delta t_2,\mathcal L_1}-p_{h_2,\Delta t_2,\mathcal L_2}\|_1
\end{align}
is a triangle inequality. It identifies useful intermediate comparators. It does not make a single mixed comparison a measurement of any one of these terms.

\section{The equal-DOF experiment and its confounders}
The early comparison used Q1 on $30\times36\times36$ cells and Q2 on $20\times24\times24$ cells. Both have 311,040 scalar DG unknowns:
\begin{equation}
 30\cdot36\cdot36\cdot8=20\cdot24\cdot24\cdot27.
\end{equation}
Equal algebraic dimension is useful for a cost comparison. Here it also changes cell size. The timesteps differed: $6.25\times10^{-4}$ for Q1 and $1.25\times10^{-3}$ for Q2. Reported covariance discrepancies of approximately 0.08558 and 0.25484 therefore did not isolate polynomial degree. Limiting, spatial resolution and temporal error could all explain the ranking.

The corrective experiment fixed the mesh and examined Q2 under controlled timestep and limiter variants. The higher-order unlimited polynomial had useful accuracy which the original corrected trajectory failed to retain. This changed the research objective from finding a satisfactory Q1 baseline to preserving the benefit of Q2 without allowing a negative density.

\scope{The same-mesh evidence invalidates the early degree-only interpretation. It does not imply that Q2 wins every accuracy-per-DOF comparison, or that a high polynomial degree compensates for unresolved spatial structures.}

\section{The limiter invariant and the roundoff bug}
The global scaling step assumes a nonnegative cell average. An additive conservation repair made some nearly empty cell averages slightly negative. A scaling factor computed from that average could then become inadmissible. The failure was a software violation of the prerequisite for the mathematical scaling argument, not a counterexample to that argument.

The repaired average stage enforces nonnegativity and the prescribed total mass together. This motivates a general rule: roundoff corrections must preserve the feasible set of the next stage. A change intended to improve one diagnostic can invalidate the hypotheses behind another operation. The current test suite checks average repair, limiter invariants and conservation separately.

\section{Attribution by removing the limiter}
Unlimited backward-Euler Q2 still produced a temporal density difference of about 0.00541, exceeding the 0.0025 gate. Thus the claim that all observed temporal error came from positivity projection was falsified. The time integrator itself needed attention.

Unlimited Q2--CN then gave a much smaller full/half-step difference, approximately $4.19\times10^{-5}$, and covariance discrepancy about 0.00371. Its measured negative probability was approximately $8.28\times10^{-4}$. The global corrected CN trajectory was admissible but its corresponding density difference was about 0.01088 and covariance discrepancy about 0.02192. The contrast supplied a controlled target: retain the accurate CN update while repairing its negative polynomial as locally as possible.

The measured negative mass is a global scalar. It does not locate where the negativity occurs or say how much polynomial alteration is required. A thin negative region can force a correction to coefficients over an entire element. Conversely, a large number of almost empty troubled cells can have little probability significance. Subsequent diagnostics therefore added correction norms and affected probability.

\chapter{From global scaling to repeated local optimisation}
\section{Terminal projection was a diagnostic, not an evolution method}
A terminal-only QP assesses how much of a completed high-order state can be recovered by local admissibility correction. It leaves intermediate states uncorrected and cannot establish stepwise positivity. Furthermore, preserving each cell's mass while making its polynomial nonnegative is infeasible if that cell's average is negative. The diagnostic exposed roughly 17,500 negative-average cells and about $4.34\times10^{-6}$ negative-average probability.

This is a direct mathematical obstruction: if $p\ge0$ almost everywhere on $K$, then $\int_Kp\,dx\ge0$. It cannot equal a strictly negative prescribed integral. The hybrid method therefore first repaired averages conservatively, then solved feasible fixed-average local QPs. Early SLSQP-based branches still recorded 13 and 14 fallback events. Those records remain failures of the zero-fallback requirement even if their terminal statistics improved.

\section{The optimiser became part of the scientific contract}
An optimiser's success flag is insufficient. The actual requirements are that the equality and inequalities are satisfied to declared tolerances, that the objective is appropriate to the physical polynomial, and that a failed local solve cannot be silently accepted. The reduced formulation eliminates one average degree of freedom, leaving a 26-dimensional QP. Cached sparse OSQP workspaces replace repeated generic dense constrained optimisation. Independent SLSQP comparisons serve as an oracle diagnostic rather than a production fallback policy.

The performance study included 2,413 local problems. Its speed gate passed, motivating repeated in-loop use. Performance certification and numerical certification remain distinct: a faster optimiser does not prove that the projection preserves the temporal order of CN.

\section{The first dynamic study and the fourth temporal level}
Repeated correction changes the numerical method. It must be tested as a complete update map
\begin{equation}
 u^{n+1}=\mathcal P_h\bigl(\mathcal T_{h,\Delta t}u^n\bigr),
\end{equation}
where $\mathcal T$ is the raw CN map and $\mathcal P$ includes average repair and polynomial QP projection. A theorem about $\mathcal T$ alone does not establish the order of this composition.

The initial three-level study gave insufficient temporal evidence. Tightening the Krylov solve and adding a fourth level produced the adjacent common-grid differences in Table~\ref{tab:temporal}. All 1,200 accepted steps were Bernstein-certified, with zero measured negative mass, optimiser failures and scalar fallbacks. Covariance discrepancies stayed near 0.00749--0.00752.
\begin{table}[htbp]\centering
\caption{Four-level local-QP Q2--CN temporal evidence. The order is computed from adjacent differences; it is not a proof of a global error expansion.}\label{tab:temporal}
\begin{tabular}{lrr}\toprule Comparison & Density $L^1$ & Observed order\\\midrule
$\Delta t,\Delta t/2$ & $5.238\times10^{-4}$ & ---\\
$\Delta t/2,\Delta t/4$ & $5.418\times10^{-4}$ & $-0.049$\\
$\Delta t/4,\Delta t/8$ & $1.620\times10^{-4}$ & $1.741$\\\bottomrule\end{tabular}
\end{table}
All three differences lie below 0.0025. The flat first pair followed by a smaller finest difference is consistent with a pre-asymptotic regime. There is only one encouraging fine-level rate estimate. Describing this as demonstrated second-order convergence would exceed the evidence.

Figure~\ref{fig:convergence} plots the temporal sequence alongside the separate mature correction study. Its panels concern different initial laws and must be interpreted separately.

\section{Identity-diffusion spatial refinement}
With the temporal configuration fixed at $\Delta t=1.5625\times10^{-4}$, the $20\to30\to45$ startup hierarchy gave differences 0.141394 and 0.0355845. Since successive linear resolution ratios are 1.5,
\begin{equation}
 p_h=\frac{\log(0.141394/0.0355845)}{\log(1.5)}\simeq3.4026.
\end{equation}
This is an observed rate for the particular common-grid observable. It exceeds neither a theorem's hypotheses nor the possibility of cancellation, projection effects or pre-asymptotic behaviour. The report correctly archives it as evidence, not formal order.

All 960 steps passed the adopted positivity and optimiser gates; the maximum mass error was about $1.20\times10^{-11}$. Covariance discrepancy decreased from 0.03033 to 0.005209 to 0.002340. The finest value lay below the Monte Carlo bootstrap p95 of about 0.00649, so that reference could no longer resolve finer covariance improvement reliably. The deterministic density differences remained informative independently of Monte Carlo sampling noise.

\chapter{Full-SPD diffusion and a stronger reference}
\section{What the mixed-diffusion gate changed}
The controlled tensor was
\begin{equation}
D=\begin{pmatrix}1&.4&.2\\.4&1&.3\\.2&.3&1\end{pmatrix},\qquad
\lambda(D)\simeq(.581212,.811321,1.607467).
\end{equation}
Its positive eigenvalues establish uniform ellipticity for this constant tensor. Off-diagonal coefficients introduce mixed derivatives. A diagonal directional finite-volume comparator cannot certify these terms by pretending that they are absent.

Because the stochastic factor satisfies $D=BB^T/2$, the configuration uses $B=\sqrt2\operatorname{chol}(D)$. Supplying the entries of $D$ in a field interpreted as $B$ would solve a different stochastic model. The controlled $30\times36\times36$ case retained the continuous Gaussian law, Q2, CN and local QP; it changed the diffusion factor explicitly.

That control gave covariance discrepancy about 0.004163, off-diagonal covariance discrepancy about 0.003806 and correlation discrepancy about 0.007067. Its mean correction was about $7.61\times10^{-5}$, maximum about $9.66\times10^{-4}$, and maximum mass error about $3.86\times10^{-12}$. This cleared the continuation gate for a spatial hierarchy.

\section{The full-SPD hierarchy}
The same $20\to30\to45$ hierarchy produced adjacent density differences 0.139500 and 0.0348401, giving an observed common-grid rate near 3.4215. Covariance discrepancies were approximately 0.02879, 0.004163 and 0.002261; the largest reported correlation discrepancies decreased from about 0.0551 to 0.00707 to 0.00178. All 960 accepted steps passed whole-cell positivity, with zero reported optimisation failures/fallbacks and mass error bounded by about $3.87\times10^{-12}$. The aggregate propagation runtime was about 4,817 seconds; 8,881,349 local QP operations were recorded.

\scope{This demonstrates robustness for one full-SPD matrix in a short Gaussian forecast. It does not certify every SPD tensor, every condition number, a rank-deficient diffusion, or a mature posterior.}

\section{Monte Carlo resolution}
The strengthened reference used one million particles and a fixed bootstrap design. The p95 covariance sampling scale was about 0.002735, closer to the finest FE discrepancy than the earlier 200,000-particle uncertainty. Common random numbers make differences between reference estimates less noisy, but they do not turn either estimate into an exact solution.

For independent samples $X_i$ and a scalar observable $g$ with finite variance,
\begin{equation}
 \operatorname{Var}\left(N^{-1}\sum_{i=1}^{N}g(X_i)\right)
 =\operatorname{Var}(g(X))/N.
\end{equation}
Thus increasing the particle count fivefold improves a standard-error scale by about $\sqrt5$, subject to the chosen observable and discretisation bias. Increasing bootstrap replicates stabilises a percentile estimate; it does not create more information about the population than the original sample contains. The reference also has an SDE timestep error, finite-horizon initialization error and histogram/smoothing choices.

The adopted comparison therefore uses MC moments and low-dimensional marginals as independent diagnostics while retaining deterministic mesh-to-mesh density differences as primary convergence evidence. The bootstrap p95 is an uncertainty scale associated with the specified resampling procedure, not a confidence guarantee for arbitrary solver bias.

\chapter{Mature densities exposed a different failure mechanism}
\section{A common physical law, an expensive representation}
The mature initial law was constructed from a Lorenz spinup ensemble and a Bayesian observation, then represented by a continuous Gaussian mixture. It is a specified smooth law with stretched lobes and low-density tails. It is neither a single Lorenz orbit nor an empirical point cloud passed directly into a finite-element density.

The original mature hierarchy retained the full-SPD tensor, Q2--CN, timestep and local-QP policy. On $20^*$, $30^*$ and $45^*$ meshes, where the asterisk abbreviates the corresponding $y,z$ counts, the initial representation corrections were roughly 0.380, 0.158 and 0.0313. These values matter before the first dynamical update. A later experiment that changes both initialization and evolution cannot isolate the latter.

The dynamic mean corrections were approximately 0.003047, 0.002789 and 0.001840. All exceed the frozen mature threshold 0.001. Affected probability stayed near 0.9986, 0.9988 and 0.9425. In contrast, covariance discrepancies improved to approximately 0.005844, 0.003636 and 0.000858; joint-TV discrepancies improved from about 0.02102 to 0.01370 to 0.005184. Accurate moments therefore coexisted with substantial full-polynomial intervention.

Adjacent density differences 0.420590 and 0.353734 gave only $p_h\simeq0.427$. All 960 baseline mature steps remained certified with zero optimiser failures/fallbacks and maximum mass error about $4.01\times10^{-12}$. The baseline aggregate gate failed on correction and density resolution, rather than loss of conservation.

\section{Deeper certificates tested a falsifiable explanation}
A Bernstein coefficient can be negative while the polynomial itself is nonnegative. Deeper subdivision could therefore conceivably rescue a large number of false-positive troubled cells. The diagnostic compared depths 4 and 8 and also searched for actual negative witnesses.

At most about 0.0788\% of cells were rescued in the tested states; correction $L^1$ reduction reached about 2.89\% in the best case, with mean about 0.779\%. On the finest terminal state, deeper analysis eliminated 39 unresolved classifications while finding additional negative witnesses. The number of witnessed negative cells changed from 130,873 to 130,881. This is evidence of genuine negative Q2 polynomials, not merely a pessimistic coefficient test.

Scalar correction changed the tested states between about 1.96 and 6.70 times as much as the QP, with mean ratio about 5.02. Greater certificate depth and more global scaling did not supply the needed accuracy gain. These findings justify retaining the local QP as the comparator; they do not prove that every possible polynomial positivity method must be worse.

\section{Average-level AFC and its recorded failures}
The AFC experiment separated initial representation from dynamics by starting candidates from the same positive projected mature state. Its low-order conservative average update accepted more than 99.9998\% of the high-order correction. Cell averages were already almost completely admissible. The subsequent polynomial correction remained close to the baseline $10^{-3}$ scale, placing the unresolved problem inside the high-order cell polynomial.

The compact scientific decision must be read alongside the wrapper status. The wrapper returned zero and marked process completion successful. The scientific classification was \code{FAILED_PREDECLARED_MATURE_BIMODAL_GATES}. Moreover, this branch recorded the failures in Table~\ref{tab:afcfail}; it cannot inherit the zero-fallback claim of the successful baseline.
\begin{table}[htbp]\centering\caption{AFC branch: reported optimizer failures and scaling fallback cells. Counts are events accumulated over the branch.}\label{tab:afcfail}
\begin{tabular}{lrrr}\toprule Mesh & Mean correction & Optimizer failures & Fallback cells\\\midrule
$20\times24\times24$ & 0.00304661 & 1,687 & 1,687\\
$30\times36\times36$ & 0.00278861 & 1,762 & 1,762\\
$45\times54\times54$ & 0.00184036 & 4,592 & 4,592\\\bottomrule\end{tabular}\end{table}
The near-unity antidiffusive acceptance addresses average positivity; it supplies no guarantee for within-cell minima. The failure counts add an independent reason to reject this particular implementation for dataset generation. The compact evidence does not by itself identify a unique numerical cause of each local optimizer failure.

\section{The DOF diagnostic and what it actually implemented}
The next diagnostic operated on Bernstein coefficients and compared a convex/simplex-type nonnegative repair with the current mass-metric QP on saved raw states. This is a useful constrained-polynomial experiment. It is narrower than constructing a conservative invariant-domain subcell flux graph for the whole anisotropic DG operator.

On the $45^*$ one-step state, its $L^1$ modification was about 0.0121874, roughly 22.10 times the QP modification. On the targeted $60^*$ state, it was about 0.00648577, roughly 40.08 times the QP modification. Neither passed the predeclared $C_{\rm DOF}<0.5C_{\rm QP}$ continuation gate. The result rejects this prototype, with its chosen coefficient constraints and metric. It leaves more elaborate flux-based subcell constructions untested.

The targeted finer-state QP intervention decreased substantially. Together with the steep reduction in initialization correction, this supported localized spatial under-resolution as the most useful next hypothesis. Refinement creates representational capacity; changing a constraint does not create missing spatial scales.

\chapter{Uniform and graded refinement: accuracy versus topology}
\section{Uniform 60 passed the original correction gate}
The $60\times72\times72$ trajectory used 311,040 cells and 8,398,080 Q2 DOFs. Its mean relative correction was $8.313\times10^{-4}$, below the original mature threshold $10^{-3}$. Positivity, mass, optimizer and statistical gates also passed. It cost about 69.4 minutes and 6.05 GiB per rank at peak.

The $45^*$--$60^*$ density difference remained about 0.20997 and the heuristic observed rate was about 0.51. The method could now handle this law's statistical and positivity requirements at the tested resolution, while full-density convergence and cost remained open. The subsequent $8\times10^{-4}$ adaptive target was a new branch gate; it must not retroactively overturn the original $10^{-3}$ criterion.

\section{Static tensor grading reproduced uniform 60}
The first graded mesh had $44\times52\times46$ cells: 105,248 cells and 2,841,696 unknowns. It concentrated fine axis spacings around the mature structures and coarsened the exterior. It reproduced the uniform-60 common-grid density to about 0.001688, reduced DOFs by 66.2\%, runtime by about 53\% and memory per rank by about 61.7\%.

Its mean correction $8.3837\times10^{-4}$ exceeded the branch's predeclared $8\times10^{-4}$ threshold. That failed aggregate decision is retained. Since it was almost the same intervention as uniform 60, the result supports efficiency of grading for reproducing that reference, without certifying the reference's spatial accuracy.

An initial marginal-statistics comparison used inconsistent histogram resolutions. Re-evaluating the matched-grid quantities left maximum marginal degradation about $1.75\times10^{-4}$, within the fixed 0.005 tolerance. This is a comparison bug correction, not a change in the solver trajectory or a post hoc relaxation of the gate.

\section{Finer grading revealed missing structure}
The $58\times68\times56$ design used 220,864 cells and 5,963,328 DOFs. Mean correction fell to approximately $1.128\times10^{-4}$ and affected probability to about 0.004179. Yet its density moved by 0.108680 relative to uniform 60 and 0.107936 relative to the previous graded mesh.

The simultaneous decrease in correction and substantial change in density is strong evidence that finer resolution reveals physical polynomial structure which the earlier meshes did not represent. A small limiter norm alone would have missed the unresolved density. The runtime rose to about 76.4 minutes despite fewer DOFs than uniform 60. Local QP activity, assembly, conditioning and solver iteration cost determine time alongside algebraic dimension.

\section{Time aggregation succeeded; tensor topology did not}
The replay used 11 representative snapshots over $0\le t\le0.05$. Reconstructed fields agreed at approximately $3.61\times10^{-11}$ in $L^1$, validating replay fidelity for those snapshots. D\"orfler-style marking identified 71 important cells at the chosen indicator fraction.

In a tensor mesh, refining intervals selected along each axis creates their Cartesian product. If selected axis intervals occupy fractions $r_x,r_y,r_z$, their product can force fine resolution in combinations never identified as important physical cells. Here the proposed design covered roughly 75\% of axis ranges and expanded to about 699,600 cells, exceeding the 311,040-cell ceiling.

The failure localized an architectural limitation. The trajectory indicator found sparse three-dimensional regions; the mesh representation could only express rectangular products of one-dimensional bands. The next implementation therefore required local hexahedral refinement without changing the PDE, polynomial degree, time integrator or projection.

\chapter{MFEM equivalence before nonconforming refinement}
\section{Porting one variable at a time}
The DOLFINx solver remained the frozen reference. MFEM supplied an alternative mesh infrastructure, with a mathematically matched Q2 DG operator on ordinary conforming hexes before enabling hanging faces. This avoids attributing a discrepancy simultaneously to a framework port and to an AMR topology.

The first structural probes found errors in physical extents and signs. A unit-length construction in an intended 60-unit direction changes Jacobians, gradients and mass, even if the number of cells agrees. Similarly, swapping the sign of a diffusion or convection contribution changes the PDE. These faults were localized in cheap operator probes and corrected before using mature trajectories as comparators.

The large-vector absolute conservation residual, approximately $-5.6\times10^{-10}$, initially appeared concerning. Normalizing the mass-functional defect across five production-scale actions gave about $1.93\times10^{-17}$. This is evidence of a near-roundoff cancellation, consistent with conservative assembly. Both absolute and normalized quantities are valuable: normalization diagnoses scaling; stepwise mass remains the practical forecast gate.

\section{Physical basis mapping}
MFEM and DOLFINx may represent the same polynomial in differently ordered local vectors. Let $T_K$ map the source coefficients to common physical nodal values. A meaningful comparison uses $T_Ku_K$ or evaluations of $p_h$, not raw vector entry equality. For an operator action, compare the mass-inverted physical derivative $M^{-1}K_hu$ after mapping, rather than mixing dual residual coefficients with primal polynomial coefficients.

The small $4\times5\times5$ probe gave mapping error about $2.55\times10^{-16}$, relative operator discrepancy $4.95\times10^{-15}$, raw CN discrepancy $2.43\times10^{-15}$, and corrected-step discrepancy $7.81\times10^{-12}$. On $30\times36\times36$, the corresponding values were about $1.19\times10^{-15}$, $7.96\times10^{-14}$, $9.43\times10^{-15}$ and $1.21\times10^{-11}$.

An early port bypassed adaptive certificate checks in one branch. That result was superseded after the missing checks were added. Passing assembly and raw CN tests cannot compensate for an omitted admissibility gate.

\section{Full uniform trajectory}
The primary clean full-equivalence artifact is
\code{runs/mfem-mature-uniform-equivalence/20260915T193445Z}.
Both frameworks used the same projected initial polynomial, 320 steps, common timestep and conservative common-grid export. The final density difference was approximately $5.26\times10^{-11}$; the covariance Frobenius difference was approximately $3.70\times10^{-11}$. Common-voxel moment differences were on the $10^{-11}$ scale. Differences of the mean and maximum correction summaries were $7.52\times10^{-15}$ and $1.72\times10^{-13}$ respectively. The saved MFEM summary does not provide every step's correction value, so this is aggregate agreement rather than a pointwise history comparison. Maximum reported mass drift was about $3.98\times10^{-13}$.

This establishes practical numerical equivalence for the tested configuration. It is stronger evidence than agreement of low-dimensional statistics alone. A later rerun had a successful solver result but a wrapper quoting error during postprocessing; that provenance problem is recorded separately. It does not supersede the clean primary run.

\section{Hanging-face gate}
Before the complete AMR trajectory, a small two-rank mesh tested representative 2:1 interfaces. Common physical weak-operator actions agreed with a conforming comparator at relative scale about $5.87\times10^{-14}$. Matching SIPG penalty geometry was essential: a framework's default inverse-normal-length convention need not equal the DOLFINx cell-diameter convention.

The NC test supports conservative, correctly scaled face assembly on those interfaces. It does not prove convergence of a long AMR forecast. MFEM's documented \code{nc_limit} mechanism can add neighbouring refinements, so actual cell counts must be measured after constructing and balancing the mesh \cite{mfemnc}.

\chapter{Local AMR, support sensitivity and the final domain sequence}
\section{First and second AMR levels}
Static AMR1 used 140,418 cells and 3,791,286 DOFs, 36.4\% fewer than the fine graded reference. It gave $D_{\rm AMR1,g_f}=0.01871569<0.02$, mean correction $2.07549\times10^{-4}$ and maximum mass error $7.53\times10^{-13}$. All recorded invariant and statistical gates passed. Propagation cost 9,780 seconds, approximately 2.13 times the fine graded runtime; peak memory was not recorded. Numerical reproduction passed while resource efficiency remained open.

AMR2 used 163,728 cells and 4,420,656 DOFs. The adjacent difference $D_{1,2}=0.0157065$ passed its 0.05 gate. Mean correction fell to $1.04980\times10^{-4}$ and mass error stayed below about $8.40\times10^{-13}$. However $D_{2,g_f}=0.0191270$ slightly exceeded $D_{1,g_f}$. One bounded adjacent difference is valuable, but it cannot establish monotone approach to an exact density.

\section{The third-depth contraction failed}
AMR3 produced $D_{2,3}=0.0106884$. The predeclared halving gate was $0.5D_{1,2}=0.00785326$. The observed ratio was about 0.6805, so the contraction gate failed even though positivity, mass, optimiser and statistical tests passed. Initial discrepancy was about 0.000812314; runtime was about 5.23 hours and peak memory about 7.16 GiB per rank.

The saved density-difference map placed 80.1\% of the AMR2--AMR3 difference outside the second-depth marking support. This gave a concrete falsifiable explanation: the replay-defined fine region was too narrow, so adding depth within it could not resolve error outside it.

\section{Support and depth were disentangled}
The expanded-support study retained existing marks, added regions covering 95\% of the final absolute density difference and one child-cell buffer, and held the finest cell size equal to AMR2. Its change from AMR3 was 0.0150712, exceeding 0.00785326. It therefore failed reproduction of AMR3, while reducing mean correction to about $2.02955\times10^{-5}$. A substantial support effect is evidence that AMR3 itself is an inadequate universal reference.

The targeted extra-depth study then held that support fixed and refined 1,889 buffered cells, covering about 39.1\% of the chosen indicator. Its difference from the support comparator was only 0.000316854 and mean correction $2.02951\times10^{-5}$. This bounded test passed its own continuation gate. Because it covered a selected fraction, it cannot certify the untested remainder by logical implication.

The broader-support study increased the discrepancy target to 99\% and retained previous expanded marks. Actual forced-neighbour refinement produced 208,787 cells rather than the estimated 204,552. Relative to the support comparator it changed the density by 0.00177608, while mean correction fell to $7.98372\times10^{-6}$. The old AMR2--AMR3 discrepancy coverage reached 99.735\%. All configured broader-support gates passed. The authoritative file is \code{broader_support_comparison.json}; inherited failed comparator files in the same directory answer older questions and retain their original classifications.

\section{Temporal sensitivity of the broader working mesh}
The half-timestep run used exactly the same initial voxel field and mesh, with 640 steps at $7.8125\times10^{-5}$. Final full/half-step density difference was $1.97036\times10^{-5}$, far below the 0.0025 gate. Mean correction per step was about $4.0187\times10^{-6}$, but the sum over 640 steps was about 0.002572, close to the 320-step sum about 0.002555. Halving a per-step intervention while doubling the number of steps is not evidence of a halved total modification.

Mass error stayed below about $2.76\times10^{-13}$ and every accepted state passed the declared gates. Runtime was approximately 4.47 hours. These data support retaining the original timestep for this working mesh and horizon; they do not supply a temporal asymptotic order from only one pair.

\section{Two aligned domain expansions completed}
The initial domain attempt stopped at the memory safeguard. Later runs released redundant mesh and assembly storage, then completed both expansions with unchanged physical interior resolution. The first expansion had 224,959 cells; the second had 242,403. They were compared on aligned conservative voxel grids, with separate added-shell probability diagnostics.

The first interior $L^1$ change was $1.32129\times10^{-12}$. The second change, independently recomputed from the two binary voxel arrays, was $1.9549638427469373\times10^{-12}$. Relative covariance changes were approximately $8.06\times10^{-14}$ and $1.59\times10^{-13}$; maximum marginal-TV changes were approximately $5.47\times10^{-13}$ and $8.17\times10^{-13}$. Measured added-shell mass was zero in the audited second comparison.

All 320 steps of the second expansion were certified, with zero reported optimiser failures/fallbacks. Maximum mass error was approximately $1.40\times10^{-12}$. Its final FE mass and export mass differed by about $7.64\times10^{-14}$. Propagation wall time was 23,804.9 seconds, including memory-pressure pauses, and peak RSS was about 6.36 GiB per rank.

\scope{The two-expansion gate is passed numerically for this mature law, tensor and $0.05$ horizon. The nearly machine-precision differences reflect dormant tested exterior effects; they are not a theorem about the unbounded-domain problem or every longer forecast. The final scientific artifact explicitly keeps production authorization false.}

\section{The resulting decision}
The late evidence supports a bounded provisional pilot using the broader-support comparator and its declared law/tensor/horizon neighbourhood. It leaves a continuum full-density error certificate open. A finite sequence of small support, timestep and domain differences does not bound an untested error component. The conditional pilot is a research decision to learn whether this target family is useful, with per-sample rejection gates; it is not permission for a large unattended production dataset.

The scientifically durable state has three separate entries: statistical and positivity gates passed; both tested domain expansions passed; mature-density continuum convergence remains open. Preserving those distinctions is the most important outcome of the complete development history.

\part{The final tested algorithm and its reproducibility}
\chapter{An executable mathematical specification}
\section{What is frozen in the tested comparator}
The current working comparator is the broader-support static MFEM NC-hex trajectory. Its mathematical ingredients are a constant full-SPD tensor, conservative Lorenz drift, reflecting total flux, Q2 DG on affine axis-aligned hexahedra, matched SIPG penalty, CN and the two-stage local-QP positivity treatment. Mesh support changed through controlled studies; the positivity mathematics remained the frozen comparator.

\begin{table}[htbp]\centering\caption{Specification of the broader mature working comparator. Counts describe the constructed mesh, including forced refinement.}\label{tab:finalspec}
\begin{tabular}{p{.32\textwidth}p{.58\textwidth}}\toprule Component & Recorded setting\\\midrule
Drift & Lorenz--63, $(\sigma,\rho,\beta)=(10,28,8/3)$\\
Diffusion & $D_{11}=D_{22}=D_{33}=1$, $D_{12}=.4$, $D_{13}=.2$, $D_{23}=.3$\\
Original box & $[-30,30]\times[-40,40]\times[-10,70]$\\
Boundary condition & $(fp-D\grad p)\cdot n=0$\\
Space & 27 local Q2 DOFs per affine hex\\
Background & $45\times54\times54$ cells\\
Static refinement & Isotropic local octree refinement, two depth levels in retained broader support\\
Actual cells / DOFs & 208,787 / 5,637,249\\
Time & $\Delta t=1.5625\times10^{-4}$, $\tau=.05$, 320 CN steps\\
Projection & Nonnegative conservative averages; fixed-average mass-metric QP\\
QP solve & OSQP, reduced 26-dimensional variables, cached workspace\\
CN linear solve & Hypre GMRES with ILU; relative tolerance $10^{-14}$, absolute $10^{-17}$, maximum 1,000 iterations, Krylov dimension 100\\
Common export & $180\times216\times216$ conservative voxel averages, float64, C ordering $(x,y,z)$\\\bottomrule\end{tabular}\end{table}

The declared settings must be read from the archived config and source revision when reproducing a specific run. A command pointing at the current source can change results if that source has changed. The table identifies the tested mathematical configuration, not a guarantee that arbitrary builds are bitwise identical.

\section{Algorithm before the first timestep}
The setup begins with a continuous normalized Gaussian-mixture law. The domain is built, selected base cells are refined, and balancing adds any required neighbours. Refinement is complete before propagation. A parallel DG space is created; the mass, transport and diffusion forms are assembled with matched geometry and quadrature.

The initial law is projected by solving
\begin{equation}
 M u^0_{\rm raw}=b^0,\qquad b_i^0=\int_\Omega p_0\phi_i\,dx.
\end{equation}
The quadrature approximates the nonpolynomial Gaussian mixture; it is not an exact symbolic Gaussian integral. Global normalization, average repair and local QP produce the accepted $u^0$. The representation correction is recorded separately from the dynamical intervention.

Construct $L=M-\Delta tK_h/2$ and $R=M+\Delta tK_h/2$ once. When these independent matrices have been assembled, the mature/AMR branch can release the original spatial matrix and redundant assembly storage. This changes storage lifetime, not the CN equation. The preconditioner is built for $L$.

\section{One complete accepted step}
\begin{lstlisting}[caption={Mathematical pseudocode for the tested two-stage update. Implementation details and abort conditions are part of the specification.}]
Given accepted positive Q2 field u and fixed matrices M, K, L, R:
  rhs := R u
  solve L raw = rhs at declared GMRES tolerances
  reject/abort if the linear solver does not converge
  compute raw cell averages, total mass and negative-average mass
  solve global weighted nonnegative-average projection
  shift each local polynomial to its repaired average
  set exactly zero-average cells to the zero polynomial
  for each local cell:
    transform into the canonical physical Q2 basis
    certify the polynomial by Bernstein subdivision
    if admissible: retain it unchanged
    otherwise solve the fixed-average mass-metric QP
    validate average, inequalities and accepted polynomial
    abort on a failed guarded projection
  reduce global mass and correction diagnostics across ranks
  accept corrected state only after its declared gates pass
  record step/time/correction; export configured snapshots
\end{lstlisting}

The algorithm's rejection behaviour distinguishes it from an optimisation call wrapped in a silent scalar fallback. The successful MFEM branch aborts when guarded local projection cannot establish acceptance. Its summary writes zero failure/fallback values on successful completion. These values are supported by the guard architecture and accepted-state checks, but they are not separately accumulated counts of recoverable failed internal attempts. A future production audit should retain explicit counters and true residual histories, even for successful runs.

\section{Conservation through all stages}
There are three mass contracts. The spatial operator satisfies $\ell^TK_h=0$, where $\ell$ represents the constant test function in the residual basis. Consequently exact CN preserves $\ell^TMu$. Stage~1 preserves its target total mass. Stage~2 preserves each repaired cell average. Therefore the composite algorithm conserves mass in exact arithmetic, provided Stage~1 is feasible and every local equality is satisfied.

The floating-point errors arise separately from assembly, inexact linear solve, scalar root finding, QP tolerances, basis transforms, reduction order and export. A final mass within $10^{-10}$ does not prove every contribution is small; cancellation is possible. The tests isolate components where practical. The MFEM summary's \code{maximum_absolute_mass_error} is maximum drift from its accepted initial mass, rather than directly from one. The two are related by $|M_n-1|\le|M_n-M_0|+|M_0-1|$. Initial masses are also reported, allowing that distinction to be checked; individual step masses are held during propagation but not all persisted in the compact summary.

No correction can repair a negative total target mass while preserving it and producing a nonnegative field. This impossible case must be rejected. Nearly zero cells also require a stable convention: dividing coefficients by a tiny average can amplify roundoff. The current projector normalizes positive-average problems and uses dedicated zero/near-vacuum handling rather than assuming every cell can be safely scaled without checks.

\chapter{Basis transformations, geometry and quadrature}
\section{Physical polynomials are the common object}
A coefficient vector has meaning only relative to basis functions, reference coordinates and cell orientation. For a canonical Q2 nodal basis on $[0,1]^3$, write
\begin{equation}
 p_K(F_K(\xi))=\sum_{i,j,k=0}^{2}c_{ijk}l_i(\xi_1)l_j(\xi_2)l_k(\xi_3),
\end{equation}
where $l_i$ are the one-dimensional Lagrange polynomials at $0,1/2,1$. The affine map is $F_K(\xi)=x_K+H_K\xi$, with diagonal $H_K$ for the tested meshes. Gradients transform by $H_K^{-T}$; volume integrals acquire $\det H_K$. Face gradients, normals and measures must follow the same map.

The physical nodal map is obtained by evaluating the framework basis at canonical points. Its inverse represents a corrected canonical polynomial back in framework coefficients. For an affine axis-aligned cell the reference mass metric is a common positive scalar multiple of the physical metric. Therefore removing the volume factor from a local QP does not change its minimizer. On curved or nonaffine cells this simplification would require revision.

\section{Quadrature has several roles}
The initialization integral is nonpolynomial because of the Gaussian mixture. The declared order 14 is numerically tested for that projection. The assembled operator uses its own integration rules: in the port, volume advection uses a specified cube rule of order 6; the combined face rule is order 14, and the isolated advection equivalence probe uses order 12. Diffusion integrators choose rules through their library conventions. It would be inaccurate to claim that every integral in both frameworks uses the same order-14 rule.

For constant diffusion, products of Q2 gradients are polynomial on affine hexes. The Lorenz drift is quadratic, so transport products have a known finite degree in each coordinate. Adequate Gaussian quadrature can integrate these polynomial forms exactly up to arithmetic error. The observed cross-framework action agreement verifies the adopted combination more directly than comparing a single nominal quadrature-order field.

\section{Penalty scaling is geometry, not a free label}
The incumbent penalty is based on $16k^2$ times a normal diffusivity over an averaged cell diameter. At $k=2$ this supplies a factor 64. MFEM's generic SIPG integrator conventions require a custom matched scaling to reproduce that geometric choice. On nonuniform cells, diameter and normal width are different. A numerical penalty constant copied from another framework need not preserve the actual physical face term.

On a hanging interface, integration is over each fine face segment. The coarse polynomial is evaluated on that same physical segment. Contributions enter the two neighbouring DG residuals with opposite transport flux signs. DG does not require the independent coarse/fine cell polynomials to agree; it requires consistent conservative face integration. Hanging-node constraints for continuous elements describe a different algebraic situation.

\chapter{Conservative export and statistical diagnostics}
\section{Why a sample image is insufficient}
A pointwise sample $p(x_V)$ can miss a narrow oscillation or bias the mass of a voxel. The common-grid field stores
\begin{equation}
 \bar p_V=|V|^{-1}\int_Vp_h(x)\,dx.
\end{equation}
For the adopted affine Q2 cells, tensor three-point Gauss integration is exact for each polynomial intersection integral. The domain and mesh are aligned so intersections have simple Cartesian geometry.

The exporter handles two current relationships. A coarse cell can contain an integer number of common voxels along every axis; it integrates separately over each voxel. A finest depth-three cell can occupy one eighth of a common voxel, with half the voxel width in every axis; its cell average contributes one eighth to the voxel average. An independent coverage array requires exactly eight eighth-coverage units for every voxel. An uncovered or multiply covered voxel triggers an error.

This implementation supports the tested aligned sizes; it does not implement arbitrary polyhedral intersections or all possible anisotropic refinements. The guard against misalignment is a valuable limitation. Silently approximating unsupported geometry would invalidate conservation and density comparisons.

\section{File semantics}
The native binary files contain only float64 values. Their dimensions, bounds, coordinate ordering, timestep and units come from companion metadata. On the baseline grid,
\begin{equation}
 h_V=(60/180,80/216,80/216),\quad
 |V|=384000/(180\cdot216\cdot216).
\end{equation}
Index $((iN_y)+j)N_z+k$ uses $z$ as the fastest varying coordinate. A consumer must check the byte count before reshaping. Endianness and float representation should be made explicit when moving files between platforms; the local evidence uses the host-written double format.

The second domain grid is $196\times232\times232$ and the first is $188\times224\times224$. Removing four outer voxels from each face of the second aligns its interior with the first. This is the alignment used in the independent recomputation. A zero-filled outer shell is a measured value at the stored resolution, not an analytic proof of compact support.

\section{Moments from FE fields and voxel fields}
Mass from voxel averages is exactly the exported integral up to arithmetic. Moments require more care. Replacing a coordinate by its voxel centre approximates $\int_Vxp_h\,dx$; it is not equal in general, even when $\bar p_V$ is exact. The DOLFINx diagnostics include native FE quadrature moments. The current MFEM mature/AMR comparison scripts compute voxel-centre moments from the conservative export, without a within-voxel variance addition. Those reported covariance numbers therefore describe that matched-grid observable, with an averaging error relative to native polynomial moments. Native FE moment export would strengthen a future audit, rather than being a capability falsely attributed to the saved MFEM summaries.

For a covariance comparison, specify whether the normalization is by $\|\Sigma_{\rm ref}\|_F$, by an entrywise denominator or by another scale. Off-diagonal entries can be small, so an entrywise relative error becomes unstable near zero. Correlation coefficients
\begin{equation}
 R_{ij}=\Sigma_{ij}/\sqrt{\Sigma_{ii}\Sigma_{jj}}
\end{equation}
are meaningful only when both variances are positive. Accurate diagonal variances alone do not validate mixed diffusion; off-diagonal covariance and correlations must also be checked.

\section{Low-dimensional statistics can miss density error}
For any bounded observable $g$,
\begin{equation}
 |\E_pg-\E_qg|\le\|g\|_\infty\|p-q\|_1.
\end{equation}
The implication runs from full-density control to bounded-observable control. A finite list of moments cannot be inverted into a general density-distance bound. There exist distinct densities with identical first and second moments. For example, symmetric mixtures with different component shapes can be adjusted to share mean zero and variance one.

Marginalization is also a contraction, so a large joint-density difference can coexist with small one-dimensional TVs. The mature experiments instantiate this issue numerically: excellent covariance and lobe probabilities appeared while adjacent full-density differences remained large.

The reported joint-TV values use a fixed coarse/smoothed Monte Carlo representation. They address agreement with that reference observable. They do not measure the exact TV between the FE field and the continuous unknown solution. The common-grid deterministic $L^1$ comparison has its own voxel-averaging scope, discussed in the convergence chapter.

\section{Final broader-support numerical statistics}
The authoritative \code{broader_support_comparison.json} reports normalized covariance discrepancy $0.00124557$ and smoothed joint TV $0.00147144$. The one-dimensional marginal TVs are approximately
\begin{equation}
 (\mathrm{TV}_x,\mathrm{TV}_y,\mathrm{TV}_z)=(0.00370487,0.00337592,0.00266795).
\end{equation}
Negative/positive-$x$ lobe probabilities are approximately $0.500000355$ and $0.499999645$; maximum lobe discrepancy against the common MC sample is approximately $0.000421645$.

The initial-projection $L^1$ change is not a dynamic statistic: the recorded total initialization correction is about $6.38820\times10^{-5}$. Dynamic mean and maximum relative corrections are $7.98372\times10^{-6}$ and $1.52993\times10^{-5}$. Maximum reported mass drift is $8.76410\times10^{-13}$. The accepted trajectory has zero recorded uncertified cells and zero successful-run failure/fallback flags. These quantities passed the configured broader-support gates, with \code{mature_density_convergence_certified=false} and \code{production_dataset_authorized=false} retained in the same report.

\chapter{Software architecture and code-to-mathematics mapping}
\section{The two implementation layers}
The Python DOLFINx layer provides model configuration, FE forecasting, analysis, local projection, validation, experiment orchestration and an existing dataset interface. The MFEM C++ branch provides the matched Q2 operator, projection, static NC-hex evolution and conservative export. Python preparation/comparison wrappers construct marking inputs and audit outputs. The separation permits cross-framework tests without claiming that every DOLFINx feature has been ported.

\begin{longtable}{p{.44\textwidth}p{.48\textwidth}}
\caption{Source map at the audited solver revision. Paths are repository-relative and indexed with hashes in the companion ledger.}\\
\toprule File or module & Mathematical/software responsibility\\\midrule\endhead
\code{lorenz_fpe/core.py} & Lorenz/noise model, DOLFINx DG assembly, forecast and Bayesian analysis, native export.\\
\code{lorenz_fpe/local_projection.py} & Average repair, Q2 mass-metric constraints, OSQP/SLSQP validation and correction diagnostics.\\
\code{lorenz_fpe/afc.py} & Controlled average-level AFC and coefficient-level diagnostic constructions.\\
\code{lorenz_fpe/finite_volume.py} & Independent diagonal-diffusion upwind finite-volume reference and SSPRK stepping.\\
\code{lorenz_fpe/validation.py} & Monte Carlo/reference quantities, mathematical and numerical diagnostics.\\
\code{lorenz_fpe/dataset.py} & Existing serial dataset writing and split/schema operations; this presence does not imply mature AMR data were generated.\\
\code{mfem/mfem_physical_equivalence.cpp} & Matched basis/operator/CN/QP implementation, static NC meshes, aligned exporter and trajectory summaries.\\
\code{mfem/dolfinx_physical_equivalence.py} & Reference physical-state and operator-action export.\\
\code{mfem/compare_physical_equivalence.py} & Mapping/action/raw-step/corrected-step gate evaluation.\\
\code{mfem/compare_mature_equivalence.py} & Full uniform cross-framework trajectory comparison.\\
\code{mfem/prepare_static_amr.py} & Indicator-to-static-base-mark construction and mixture serialization.\\
\code{mfem/prepare_support_test.py} & Density-discrepancy support expansion.\\
\code{mfem/prepare_support_depth.py} & Bounded additional depth on frozen support.\\
\code{mfem/run_broader_support.py} & Broader-support pipeline, provenance and comparison orchestration.\\
\code{mfem/run_domain_sensitivity.py} & Aligned padded domains with fixed interior physical resolution.\\
\code{mfem/memory_backpressure.py} & Process-group pause/resume under available-memory thresholds.\\
\code{scripts/check} & Static checks, compilation, tests and whitespace checks.\\
\code{scripts/verify-math} & Executable invariants and mathematical evidence consistency.\\
\code{scripts/build_solver_monograph_evidence.py} & This book's hashed evidence index, direct array checks and actual-data figures.\\\bottomrule
\end{longtable}

\section{Provenance accompanies the numerical payload}
An immutable run should retain config, source revision, dirty state, command, seeds, environment, status, comparison diagnostics and input hashes. The run wrappers record many of these pieces. Completeness varies among older runs; a missing field must be labelled missing rather than reconstructed as though it were measured.

The compact tracked root reports are convenient summaries. Immutable run decisions are the primary source for a particular experiment. Inherited comparison files must be interpreted according to their named question. A \code{status.json} can describe process state; a \code{scientific_decision.json} or named comparison classifies the scientific gates. Keeping both prevents process completion from becoming accidental method certification.

\section{The verifier's scope}
The verifier checks algebraic invariants, diffusion assumptions, limiter consistency and agreement between compact claims and evidence. It is an executable audit of declared contracts. Passing it does not make an open research gate disappear. In particular a field such as \code{method_certified: false} can coexist with \code{VERIFY-MATH: PASS}: the latter means that the records consistently disclose the former.

Regression tests protect bugs already discovered. They cannot exhaust every posterior, tensor, nonconforming topology or machine-memory condition. A production dataset needs sample-level diagnostics and explicit rejection rules in addition to repository tests.

\chapter{Runtime, memory and recoverable local execution}
\section{Cost must be measured at the right boundary}
Let $T_{\rm total}$ contain preparation, assembly, initialization, propagation, export and comparison. A timer started at the first timestep measures only a subset. The MFEM trajectory's recorded propagation wall clock includes time spent stopped by the memory-pressure guard. Peak per-rank RSS also differs from aggregate machine use: ranks can share mapped pages, while Python wrappers and the desktop consume additional memory.

\begin{table}[htbp]\centering\caption{Rounded mature-trajectory cost evidence. Timers and RSS coverage differ; entries are descriptive rather than standardized benchmarks.}\label{tab:cost}
\begin{tabular}{lrrr}\toprule Trajectory & Q2 DOFs & Propagation & Peak GiB/rank\\\midrule
Uniform $60^*$ & 8,398,080 & 69.4 min & 6.05\\
First graded & 2,841,696 & 32.6 min & 2.32\\
Fine graded & 5,963,328 & 76.4 min & about 4.45\\
AMR1 & 3,791,286 & 2.72 h & unrecorded\\
AMR2 & 4,420,656 & about 2.58 h & 6.33\\
AMR3 & variable refined topology & 5.23 h & 7.16\\
Broader support & 5,637,249 & 2.93 h & 6.98\\
Half timestep & 5,637,249 & 4.47 h & 6.80\\
Second domain & 6,544,881 & 6.61 h incl. pauses & 6.36\\\bottomrule\end{tabular}\end{table}

The fine-graded memory estimate here is rounded from its reported reduction relative to uniform 60; it is not an independently sampled peak. The later AMR values come from saved wrapper/resource records. Fewer DOFs can still mean more expensive faces, poorer conditioning, more QPs or a less favourable parallel partition. Static adaptation should be evaluated on accuracy per minute and memory as well as cell count.

\section{Pause/resume preserves a live process}
The guard reads Linux \code{MemAvailable}, a kernel estimate of memory available for starting applications without swapping \cite{linuxproc}. It sends a stop signal to the child process group when free headroom crosses the floor, then resumes that group when headroom exceeds a higher threshold. The completed run used a 2 GiB pause / 4 GiB resume hysteresis. The source defaults later moved to 2/2 GiB; already-running wrappers retain values loaded when launched.

Stopping a process retains its address space. This mechanism is neither checkpoint/restart nor disk-backed out-of-core computation. If the paused solver itself holds most of the memory, external recovery may be needed before it can resume. A polling guard can miss a rapid allocation spike and cannot guarantee prevention of the kernel OOM killer. Zero in a recorded minimum-available-memory field is not sufficient evidence of a physically accurate minimum; the audit treats that monitor datum cautiously.

Low CPU with high resident memory can be consistent with a paused process. Progress checks should examine fresh step logs, pause state and service status rather than process existence alone. A saved final array and completed comparison are much stronger evidence of completion than an idle process.

\section{Persistent jobs and estimates}
The project uses persistent local job services to survive an interactive chat closing. An ETA extrapolates completed steps only after representative progress exists:
\begin{equation}
 \widehat T_{\rm remaining}=(N-n)\frac{T_n-T_{n_0}}{n-n_0}+T_{\rm export/comparison}.
\end{equation}
QP activity and pauses can make this inaccurate. A useful estimate reports a range and identifies whether paused time is included. This monograph records completed measured durations; it launches no new PDE trajectory.

\chapter{Reproducing the evidence and planning a bounded pilot}
\section{A source repository is only one part of reproduction}
The public Git repository contains code and compact evidence reports. Large ignored arrays, archived particles, mature-mixture parameters, replay indicators and local compiled libraries are separate inputs. The AMR wrappers currently refer to local absolute environment prefixes and timestamped run paths. Cloning the repository alone cannot recreate every archived experiment until these external inputs are supplied or regenerated with their recorded seeds and settings.

The following sequence describes a controlled reproduction rather than an assurance of one-command portability:
\begin{enumerate}
\item Check the source revision and local working tree; preserve any unrelated edits.
\item Read \code{environment.yml} and the lock file; create the declared DOLFINx/PETSc/OSQP environment and an MFEM/Hypre MPI build compatible with the C++ branch.
\item Run \code{./scripts/check} and \code{./scripts/verify-math}; retain their full outputs.
\item Locate or regenerate the mature mixture, MC particles, replay indicator and comparator arrays. Verify dimensions, seeds and hashes against the run records.
\item Run physical-basis equivalence before interpreting a changed compiler/framework build as the same method.
\item Prepare a new timestamped mesh design and immutable run directory; retain original data unchanged.
\item Launch the selected wrapper persistently with explicit ranks and memory thresholds. Record the actual command and package/build versions.
\item Check accepted-state diagnostics and the named scientific comparison. Do not substitute wrapper return code for scientific classification.
\end{enumerate}

The earlier shell wrappers compile using the local MFEM MPI compiler, include/link paths for MFEM, Hypre and OSQP, and an rpath for those local libraries. These are machine-specific build contracts. A portable release should replace absolute prefixes by configured dependencies, archive reference inputs with checksums and specify binary format metadata.

\section{Dependency versions and MPI separation}
The pinned Python project environment and the MFEM environment are distinct. The Python declaration records Python 3.12.13, DOLFINx/Basix 0.11.0, UFL 2026.1.0, PETSc/petsc4py 3.25.4, NumPy 2.5.2, SciPy 1.18.0, OSQP 1.1.3 and Open MPI 5.0.10. The local MFEM environment inspected during this audit contains MFEM 4.8, Hypre 2.32.0, Open MPI 5.0.7 and GCC/G++ 15.3.0. The archived broader/domain environment output independently confirms its MPI launcher version 5.0.7. A current package inventory cannot retroactively establish every older build's complete dependency set.

Do not substitute the Python environment's MPI launcher for the MFEM environment's launcher without verifying binary compatibility. The wrappers explicitly select their launcher and linked libraries. They set \code{OMP_NUM_THREADS=1} and \code{OPENBLAS_NUM_THREADS=1} to avoid uncontrolled thread multiplication inside eight MPI ranks. These settings describe the tested execution design rather than an established optimal rank count.

\section{Important experiment entry points}
The following commands create new runs and may consume substantial local resources. They are documented for reproduction and were not executed during preparation of this book. Each requires the referenced inputs and environment. Historical configurations retain their predeclared gates.
\begin{lstlisting}
# DOLFINx startup temporal, spatial and mixed-diffusion studies
./scripts/run-experiment experiments/local-q2-cn-tight-ksp-refinement.yaml
./scripts/run-experiment experiments/local-q2-cn-spatial-refinement.yaml
./scripts/run-experiment experiments/local-q2-cn-full-spd-spatial-refinement.yaml
# Cross-framework gates before interpreting a changed MFEM build
bash mfem/run_physical_equivalence_gate.sh
bash mfem/run_mature_uniform_equivalence.sh
# Static AMR preparation/evolution uses archived replay and mixture inputs
bash mfem/run_static_amr_reproduction.sh
./scripts/run-in-env python mfem/run_broader_support.py
./scripts/run-in-env python mfem/run_broader_support.py --half-timestep
./scripts/run-in-env python mfem/run_domain_sensitivity.py --help
\end{lstlisting}
The broader-support wrapper expects an already compiled selected executable; the static-AMR wrapper contains its build command. Domain restart options validate and reuse a completed first expansion instead of repeating it. A new reproduction should retain the exact selected executable hash, not merely the C++ source name.

\section{Reproducing this book}
The editable source is \code{docs/solver_monograph/book.tex} with included chapter files. The evidence generator writes the JSON ledger, appendices and figure PDFs under \code{output/pdf/solver_monograph}. The PDF is built from the source directory with a standard LaTeX installation:
\begin{lstlisting}
./scripts/run-in-env python scripts/build_solver_monograph_evidence.py
cd docs/solver_monograph
pdflatex -interaction=nonstopmode -halt-on-error \
  -jobname=lorenz_solver_monograph -output-directory=../../output/pdf book.tex
pdflatex -interaction=nonstopmode -halt-on-error \
  -jobname=lorenz_solver_monograph -output-directory=../../output/pdf book.tex
\end{lstlisting}
The generated source ledger describes the state at generation. Rebuilding after a code change creates a new evidence snapshot and must not be misrepresented as identical provenance. The accompanying manifest records the delivered PDF hash and validation.

\section{What a pilot candidate could contain}
The discussed pilot is a small research dataset, with tens to low hundreds of trajectories, not a completed archive. An illustrative planning size was 48 trajectories with conservative coarsening to a $60\times72\times72$ learning grid and a three-dimensional Fourier-neural-operator candidate. No trained neural operator, achieved generalization result or successful pilot is evidenced by that proposal.

For each sample, the input contract should include initial density, diffusion tensor or noise factor with its convention, forecast horizon, domain/grid metadata and any observation-conditioning information. The target is the conservatively exported forecast density. Conditioned posteriors should be separate from unconditioned laws in both labels and held-out partitions. A split by trajectory family, initial law and diffusion matrix avoids near-duplicate train/test leakage.

A sample should be rejected for failed accepted-state positivity, mass, linear residual, optimizer, export coverage or boundary diagnostics. Record both intervention norms and affected probability. Testing whether these predict downstream neural-operator error is scientifically useful, but correlation would not establish that the limiter causes that error; under-resolution can cause both quantities.

\section{Remaining mathematical and scientific uncertainty}
The recorded evidence supports the selected mature configuration strongly at low-dimensional statistics, positivity, conservation, bounded support/depth/timestep sensitivity and two domain expansions. It does not prove the full-density continuum error is below a specified tolerance. The gate sequence includes failed contraction tests, and voxel averaging can hide subvoxel differences.

The untested extensions include other mature posteriors, longer horizons, highly anisotropic or ill-conditioned SPD tensors, rank-deficient noise, moving mesh adaptation and robust production-scale storage/restart. A pilot may investigate these limits with controlled rejection. Large-scale authorization requires a separate explicit project decision based on pilot quality and held-out numerical evidence.

The method's mathematical strength is an understandable conservative structure plus a local admissibility optimisation whose hypotheses can be checked. Its scientific strength is a history of falsifiable diagnostics. The limitations remain visible so that future changes can improve the solver without inheriting unsupported claims.

\appendix
\chapter{Local polynomial algebra and constrained optimisation}\label{app:local-algebra}
\section{Deriving the Q2 mass metric}
On $[0,1]$, use nodal Lagrange functions at $0,1/2,1$. Direct polynomial integration gives
\begin{equation}
 H_1=\left(\int_0^1l_i l_j\,d\xi\right)_{i,j=0}^{2}
 =\frac1{30}\begin{pmatrix}4&2&-1\\2&16&2\\-1&2&4\end{pmatrix},
 \qquad a_1=\begin{pmatrix}1/6\\2/3\\1/6\end{pmatrix}.
\end{equation}
For any nonzero coefficient vector $c$, $c^TH_1c=\int_0^1(\sum c_il_i)^2>0$, because a nonzero polynomial cannot vanish almost everywhere. This proves positive definiteness without inspecting eigenvalues. The negative off-diagonal entries are compatible with a positive quadratic form; they also explain why positivity of nodal coefficients and positivity of the function are different questions.

On the cube,
\begin{equation}
 H=H_1\otimes H_1\otimes H_1,\qquad
 a=a_1\otimes a_1\otimes a_1.
\end{equation}
Fubini's theorem applies to continuous polynomial products on a compact cube. On an affine cell, the physical mass matrix is $|K|H$. Therefore
\begin{equation}
 \|p_c-p_{c_0}\|_{L^2(K)}^2=|K|(c-c_0)^TH(c-c_0).
\end{equation}
This is the physical interpretation of the QP objective; choosing an arbitrary Euclidean norm of nodal values would weight the polynomial differently.

\section{Bernstein conversion and subdivision}
For degree two, the Bernstein basis is $(1-\xi)^2$, $2\xi(1-\xi)$, $\xi^2$. Evaluating at the three nodal locations gives
\begin{equation}
 \beta_0=c_0,\quad \beta_2=c_2,\quad
 \beta_1=2c_1-\tfrac12(c_0+c_2),\qquad
 T=\begin{pmatrix}1&0&0\\-1/2&2&-1/2\\0&0&1\end{pmatrix}.
\end{equation}
The cube conversion is $T\otimes T\otimes T$, with ordering matched to the chosen canonical vector. A midpoint split uses de Casteljau averages:
\begin{align}
 \beta^{L}&=(\beta_0,(\beta_0+\beta_1)/2,(\beta_0+2\beta_1+\beta_2)/4),\\
 \beta^{R}&=((\beta_0+2\beta_1+\beta_2)/4,(\beta_1+\beta_2)/2,\beta_2).
\end{align}
Every new coefficient is a convex combination of previous coefficients. Subdivision preserves the exact polynomial while making local bounds tighter. Tensor subdivision creates eight subcubes; their coefficient inequalities define the fixed-subcell constraint matrix $C$ used by the local QP.

All Bernstein basis functions integrate to $1/3$ in one dimension and $1/27$ in three dimensions. Consequently the cell average is the mean of whole-cell Bernstein coefficients, even though it is a weighted mean of Lagrange coefficients. Confusing these two statements changes the conservation constraint.

\section{Nullspace elimination of the equality}
Let $a\ne0$ and choose $Z\in\R^{27\times26}$ whose columns span $\ker a^T$. Every vector with average $m$ is representable as
\begin{equation}
 c=m\mathbf1+Zy,
\end{equation}
because $a^T\mathbf1=1$. The reduced objective is
\begin{equation}
 \tfrac12y^T(Z^THZ)y+y^TZ^TH(m\mathbf1-c_0)+\text{constant},
\end{equation}
with inequalities $CZy\ge-Cm\mathbf1$ (or the declared interior margin). The Hessian is SPD: $y\ne0$ implies $Zy\ne0$, so $y^TZ^THZy>0$. The equality is now satisfied algebraically, subject only to the numerical accuracy of $Z$ and transformations.

For $m>0$, normalizing by $m$ makes the problem's scale comparable across probability-rich and low-density cells. The returned vector must still be checked in physical scale. A fixed absolute tolerance acceptable near a peak may overwhelm the mass of a tail cell; normalization and explicit tolerances are both necessary.

\section{KKT conditions as an independent check}
For the full problem with $Cc\ge b$ and $a^Tc=m$, let $\mu\ge0$ multiply $b-Cc$ and $\lambda$ multiply the equality. At an optimum,
\begin{align}
 H(c-c_0)+\lambda a-C^T\mu&=0,\\
 a^Tc=m,\quad Cc-b&\ge0,\quad\mu\ge0,\\
 \mu_i(Cc-b)_i&=0.
\end{align}
These are primal feasibility, dual feasibility, stationarity and complementarity. For a strictly convex feasible quadratic problem, any point satisfying them is the unique minimizer. In finite precision, scaled residuals of these relations are diagnostic. The operational acceptance still checks the returned polynomial and mass; an optimiser's termination message cannot replace those checks.

\section{Objective guarantees and their limits}
If the input is feasible, its zero correction is uniquely optimal. If a scalar-scaled polynomial is feasible, the QP objective cannot exceed its objective. For two coefficient norms related by constants $\alpha,\beta>0$, bounds can be transferred with those constants; there is no universal identity between mass-metric and $L^1$ correction. Nor does minimal distance to the feasible set imply minimal distance to the unknown PDE solution. The raw target itself can be inaccurate.

\chapter{PDE spaces, stability and linear algebra}
\section{Why a weak formulation uses different regularity}
$L^2(\Omega)$ consists of square-integrable functions, identified up to equality almost everywhere. $H^1(\Omega)$ consists of functions in $L^2$ whose weak first derivatives are also in $L^2$. A weak derivative is defined by integration against smooth compactly supported tests, avoiding pointwise differentiability assumptions. The conforming diffusion weak form naturally uses $H^1$ because it contains one derivative of each function.

DG functions lie in broken $H^1$: they are $H^1$ on each cell but can jump at faces. Ordinary distributional derivatives across jumps include face distributions, so simply inserting a discontinuous function into a continuous $H^1$ formula is incomplete. Numerical flux, consistency and penalty terms supply the missing interface information. Trace inequalities control face quantities in terms of cell norms and gradients; inverse inequalities bring explicit powers of polynomial degree and cell size.

SIPG coercivity requires a sufficiently large penalty relative to those trace constants, shape regularity and tensor bounds. The project's chosen penalty is numerically verified for its adopted geometry. This monograph derives the usual absorption argument in the diffusion chapter, without pretending that one numerical constant is valid for every mesh and tensor.

\section{Scalar stability functions explain the timestep change}
For $y'=\lambda y$, the amplification factors are
\begin{equation}
 R_{\rm FE}(z)=1+z,\quad R_{\rm BE}(z)=\frac1{1-z},\quad
 R_{\rm CN}(z)=\frac{1+z/2}{1-z/2},\qquad z=\Delta t\lambda.
\end{equation}
If $\Re z\le0$, then $|1+z/2|\le|1-z/2|$, so CN is A-stable on this scalar test equation. Yet $R_{\rm CN}(z)\to-1$ as real $z\to-\infty$: stiff modes alternate rather than being strongly damped. BE tends to zero and is L-stable. This explains why CN can improve smooth temporal accuracy while leaving oscillatory admissibility problems. Scalar A-stability does not prove nonnegative coefficients, nonnegative polynomials or contraction for a nonsymmetric nonnormal matrix in an arbitrary norm.

Taylor expansion of a smooth exact solution about the midpoint gives
\begin{equation}
 \frac{y(t+\Delta t)-y(t)}{\Delta t}
 -\tfrac12\{y'(t+\Delta t)+y'(t)\}
 =-\frac{\Delta t^2}{12}y'''(t)+O(\Delta t^3).
\end{equation}
The unscaled step defect is $O(\Delta t^3)$. Under a uniform finite-time stability bound and sufficient smoothness this yields global $O(\Delta t^2)$ error for raw CN. Repeated nonlinear projection introduces additional defects; a separate bound is required to preserve that order. The empirical four-level test addresses this composition, rather than importing the raw-CN proof unchanged.

\section{Residual error and preconditioning}
For $Lx=b$ and approximate $\widetilde x$, residual $r=b-L\widetilde x$ satisfies
\begin{equation}
 x-\widetilde x=L^{-1}r,\qquad
 \|x-\widetilde x\|\le\|L^{-1}\|\,\|r\|.
\end{equation}
A small relative residual can coexist with a larger forward error if conditioning is poor. A preconditioned residual also depends on the chosen preconditioner. GMRES minimizes a residual over a Krylov space for nonsymmetric systems; its restarted variant retains a bounded basis at the cost of potentially slower convergence. ILU approximates factorization sparsely. Neither name establishes the achieved physical-field error; production-scale basis-transformed probes provide additional evidence.

The mass matrix is SPD and block-local for DG, making PCG a sensible initial-projection solver. The implementation uses Hypre diagonal scaling, relative $10^{-14}$, absolute $10^{-16}$ and maximum 500 iterations for that solve. The CN matrix contains advection and is nonsymmetric, so the algorithm uses GMRES rather than importing the SPD assumptions of PCG.

\section{Dimensional complexity}
For a tensor grid with $n$ cells per axis in dimension $d$, the number of Qk DG unknowns scales as $n^d(k+1)^d$. Doubling every spatial resolution multiplies cells by $2^d$. At $d=3$ this is eightfold; at $d=6$ it is sixty-four-fold. This dimensional growth explains why resolving a three-dimensional law can already be expensive and why direct high-dimensional density propagation needs fundamentally different representations.

Element-local correction can be parallelized, but global mass repair requires reductions and sparse implicit solves require communication. A graph partition with many hanging faces can increase communication and face work. Peak sparse-matrix storage depends on stencil connectivity and fill, so DOF count alone cannot predict memory. Low-dimensional latent representations, particles and neural operators can be useful future tools, but their approximation errors need separate validation.

\chapter{Verification, gates and assumptions}\label{app:assumptions}
\section{Verification and validation answer different questions}
Verification asks whether the code solves its specified discrete problem correctly. Unit tests of average repair, basis maps, operator conservation and conservative export are examples. Validation asks whether the specified model and numerical solution adequately represent the intended physical/statistical problem. Full-SPD reference comparisons, mature-state refinement and domain sensitivity address parts of validation. Neither substitutes for the other.

\begin{longtable}{p{.27\textwidth}p{.31\textwidth}p{.33\textwidth}}
\caption{Representative predeclared gates and their evidenced outcomes. Thresholds belong to their named experiments.}\\
\toprule Gate & Definition / threshold & Outcome and scope\\\midrule\endhead
Accepted-state positivity & Whole-cell Bernstein certification at each step; zero measured negative mass & Passed on final mature comparator and both domain expansions. Floating acceptance uses declared tolerances.\\
Mass & MFEM reported maximum $|M_n-M_0|<10^{-10}$, with initial normalization error also recorded & Final second domain drift approximately $1.40\times10^{-12}$; absolute-to-one bound adds $|M_0-1|$.\\
Optimizer & Zero failures and fallbacks & Successful guarded MFEM trajectories pass; AFC hierarchy failed with thousands of events.\\
Startup temporal & Every adjacent exported $L^1<0.0025$ & Four-level sequence passes; clean second-order theorem absent.\\
Startup spatial & Decreasing adjacent differences; positive observed rate & Identity and full-SPD hierarchies pass; rates 3.4026 and 3.4215 are observed common-grid diagnostics.\\
Original mature correction & Mean relative intervention $\le10^{-3}$ & Uniform $60^*$ passes; earlier $20^*$--$45^*$ fail.\\
Graded branch correction & Mean relative intervention $<8\times10^{-4}$ & First graded fails at $8.3837\times10^{-4}$; finer/broader pass.\\
DOF prototype & $C_{\rm DOF}<0.5C_{\rm QP}$ & Failed: ratios about 22.10 and 40.08 on one-step states.\\
AMR1 reproduction & $D_{{\rm AMR1},g_f}<0.02$ & Passed at 0.01871569; resource efficiency open.\\
AMR2 continuation & $D_{1,2}<0.05$ & Passed at 0.0157065; one pair supplies no asymptotic proof.\\
AMR3 contraction & $D_{2,3}<0.5D_{1,2}$ & Failed: 0.0106884 exceeds 0.00785326.\\
Broader temporal & Final full/half $L^1<0.0025$ & Passed at $1.97036\times10^{-5}$.\\
Domain density & Interior exported $L^1<0.0025$, initial interior $L^1<10^{-6}$ & Both expansions pass at near-roundoff differences.\\
Domain statistics & Relative covariance and marginal TV changes $<0.005$; absolute mean changes $<0.005$; principal direction angles $<1$ degree & Both aligned expansions pass numerically.\\
Domain boundary & Outer-layer probability $<10^{-6}$; mass/export error $<10^{-10}$; negative mass $\le10^{-13}$ & Recorded shell/layer probability is negligible, with certified accepted states.\\
Production & Explicit authorization after required evidence & False in final scientific artifact.\\\bottomrule\end{longtable}

\section{Repository assumptions translated into mathematics}
\begin{description}
\item[A1] Affine axis-aligned hexahedra support the tensor-basis transforms and exact polynomial voxel integration. The file's original Q1 implication is historical; Q2 uses Bernstein constraints instead of vertex positivity alone.
\item[A2] Zero total probability flux is the finite-domain boundary target. It is distinct from separately imposing zero drift flux and zero diffusive flux.
\item[A3] Default $B=I$ gives $D=I/2$. The production target includes constant general $B$ and full-SPD $D$; a diagonal-only reference has narrower applicability.
\item[A4] Rank deficiency does not establish uniform ellipticity. Hypoelliptic regularity, if relevant, requires a separate analysis of drift/noise brackets.
\item[A5] Backward Euler is a validated baseline; it is no mandate to keep its first-order temporal error when CN is better supported.
\item[A6] Adopted projection quadrature is numerically adequate for tested laws. This does not establish exact Gaussian integration or universal adequacy for narrower future mixtures.
\item[A7] Nonnegative fixed-subcell Bernstein coefficients are sufficient and can be conservative. They are not a necessary condition for every positive Q2 polynomial.
\item[A8] Q1 is not certified for large dataset generation. Positivity alone cannot establish target accuracy.
\item[A9] Exact-native-law MC sampling is an independent short-time reference subject to finite samples and SDE discretisation. It is not an exact high-resolution density oracle.
\item[A10] A current method ranking is revisable when controlled evidence changes accuracy, cost or applicability.
\item[A11] DOLFINx is a reference implementation rather than a permanent framework restriction. A port must meet model and equivalence contracts.
\end{description}
These assumptions are quoted conceptually from the current registry. Their status matters as much as their wording: several are deliberately recorded as false propositions to prevent recurring mistakes.

\section{A failure catalogue with bounded conclusions}
\begin{longtable}{p{.28\textwidth}p{.31\textwidth}p{.31\textwidth}}
\caption{What each rejected path established, and what it left untested.}\\
\toprule Episode & Established failure & Boundary of inference\\\midrule\endhead
Equal-DOF ranking & Degree, mesh and dt confounded & Does not rank all Q1/Q2 cost tradeoffs.\\
Roundoff average fix & Broke nonnegative-average prerequisite & Fix concerns implementation, not a failure of the scaling proof.\\
Unlimited BE refinement & Temporal density gate failed without limiter & Projection cannot be the sole explanation.\\
Global scaling CN & Accurate shape strongly degraded & Other local feasible methods remain possible.\\
Pure fixed-mass QP & Negative-average cells infeasible & Average repair is mathematically required for that constraint.\\
Three temporal levels & Insufficient finest decrease & Fourth level supplies new numerical evidence, not a theorem.\\
Mature coarse Q2 & Large initialization/dynamic corrections and slow density contraction & Accurate moments alone do not certify density.\\
Deeper certificates & Almost no relief; negative witnesses persisted & A certificate cannot repair a negative polynomial.\\
Average-level AFC & Averages already admissible; polynomial intervention and fallbacks remained & No rejection of every subcell flux construction.\\
Bernstein-DOF prototype & Much larger $L^1$ change than QP & Euclidean repair and strict constraints differ from physical-metric subcell QP.\\
First grading & Failed $8\times10^{-4}$ gate by a small margin & Efficient reproduction of uniform 60 still supported.\\
Tensor replay mesh & Cartesian product exceeded cell ceiling & Time-aggregated indicator itself remained useful.\\
MFEM early probes & Extent, signs, penalty/certificate issues & Corrected controlled port later passed equivalence.\\
AMR3 & Halving gate failed & Invariant/statistical successes remain valid components.\\
Narrow support & Density change outside marked region dominated & Depth alone did not address support error.\\
Initial domain job & RAM safeguard stopped completion & No scientific boundary conclusion from incomplete trajectory.\\\bottomrule\end{longtable}

\chapter{Reading figures and further study}
\section{A conceptual hanging-face diagram}
\begin{figure}[htbp]\centering
\begin{tikzpicture}[scale=.9]
\draw[thick] (0,0) rectangle (3,3);\draw[thick] (3,0) rectangle (6,3);
\draw (3,1.5)--(6,1.5);\draw (4.5,0)--(4.5,3);
\draw[very thick,blue] (3,0)--(3,3);
\node at (1.5,1.5) {coarse DG cell};
\node[align=center] at (4.5,3.55) {refined neighbouring cells};
\draw[-{Stealth},red,thick] (2.4,.75)--(3.6,.75);
\draw[-{Stealth},red,thick] (2.4,2.25)--(3.6,2.25);
\node[align=center] at (3,-.65) {shared physical flux on each fine face segment};
\end{tikzpicture}
\caption{Conceptual two-dimensional section of a nonconforming hex interface. This schematic is not a rendering of the archived 3-D mesh. Independent cell polynomials exchange probability through integrated shared face segments.}\end{figure}

\section{Which requested figures are supported}
The density and marginal figures use archived full three-dimensional voxel arrays. The convergence figure plots reported comparisons. The dependency diagram and hanging-face illustration are conceptual. An attractor trajectory or stochastic path figure would need an appropriately archived path with a documented horizon and law. The present short forecast arrays do not supply that evidence. No synthetic butterfly or invented trajectory was inserted to make the figures resemble a familiar Lorenz image.

The initial and final density panels describe only $0\le t\le0.05$ evolution of the specified mature posterior. They neither illustrate convergence to a stationary law nor establish that the deterministic strange attractor is the support of the noisy density. Full-rank noise changes that interpretation.

\section{Worked diagnostic exercises}
\begin{enumerate}
\item For $q(\xi)=(\xi-1/4)^2-1/100$, compute its three Q2 nodal values, Bernstein coefficients and minimum. Determine which tests detect the negative region. Then integrate $q$ to show that positive average does not imply positive polynomial.
\item Verify $\mathbf1^TH_1\mathbf1=1$ and $H_1\mathbf1=a_1$. Tensorize the identity to show that the Q2 constant's mass is the cell volume.
\item Take equal-volume raw averages $(-.1,.4,.7)$. Solve the water-filling projection with target sum one. The active positive components lose the same amount; check that the resulting vector is $(0,.35,.65)$ and derive its multiplier.
\item On unequal volumes $(1,2,1)$ with the same averages, recompute the weighted projection. Explain why using the previous unweighted answer would conserve the wrong mass.
\item Construct two symmetric normalized densities with mean zero and variance one but different shapes. Use this to refute inference of full-density closeness from covariance alone.
\item With $D_1=.0157065$ and $D_2=.0106884$, compute the contraction ratio and compare it with a halving gate. Explain why no nonuniform formal order follows from this ratio alone.
\item Let $p-q$ alternate sign within every export voxel with zero voxel integral. Verify $P_V(p-q)=0$ while $\|p-q\|_1>0$. This is the edge case behind the export-contraction qualification.
\item For a stiff mode $z=-100$, compare BE and CN amplification. Relate the signs to positivity and the magnitudes to damping, keeping those two questions separate.
\end{enumerate}
The exercises make the method's failure modes reproducible with small calculations. None requires accepting an undocumented numerical run as a premise.

\chapter{Complete recorded Git chronology}
\label{app:git}
Commit dates record when changes were logged. They do not prove when uncommitted experiments ran; immutable run timestamps provide the separate execution chronology.
\small
\begin{longtable}{p{.12\textwidth}p{.20\textwidth}p{.56\textwidth}}
\toprule Commit & Recorded date & Purpose\\\midrule\endhead
127f8ac & 2026-09-09 & chore: establish numerical research baseline\\
627ca4d & 2026-09-09 & chore: add agent-legible numerical research harness\\
1a7dc37 & 2026-09-09 & docs: keep numerical research scaffold adaptive\\
5d7973a & 2026-09-09 & docs: define minimal change logging policy\\
d0b2668 & 2026-09-10 & Audit Lorenz Fokker-Planck method selection\\
242ccf6 & 2026-09-10 & Implement same-mesh Q2 positivity experiment\\
ad9258d & 2026-09-10 & Preserve Q2 states for timestep validation\\
3d8e3b2 & 2026-09-11 & Record same-mesh Q2 gate results\\
43b83d2 & 2026-09-11 & Repair positivity limiter invariants\\
ef6ea7d & 2026-09-11 & Record corrected and unlimited Q2 evidence\\
2ffda91 & 2026-09-11 & Record Crank-Nicolson Q2 diagnostic\\
0492599 & 2026-09-11 & Reject global scaling for Crank-Nicolson Q2\\
7df38b7 & 2026-09-11 & Synchronize Q2 Crank-Nicolson evidence\\
561ec4e & 2026-09-11 & Add local Q2 positivity projection diagnostic\\
137e101 & 2026-09-11 & Handle zero-average local projection cells\\
aeaeac8 & 2026-09-11 & Certify local projection against roundoff\\
a27bc56 & 2026-09-11 & Advance local Q2 projection to dynamic test\\
03d7803 & 2026-09-11 & Profile dynamic local Q2 projection\\
7f5e97c & 2026-09-11 & Specialize local Q2 positivity projection\\
83978df & 2026-09-11 & Record dynamic local Q2 projection decision\\
7226fdc & 2026-09-12 & Predeclare tighter local Q2 refinement\\
6ae6b1e & 2026-09-12 & Advance local Q2 CN to spatial refinement\\
c8d382b & 2026-09-12 & Refresh certified evidence manifest\\
c573c7a & 2026-09-13 & Advance certified solver to full-SPD gate\\
4a5336a & 2026-09-13 & Refresh full-SPD evidence manifest\\
619b8b4 & 2026-09-13 & Archive passing full-SPD control\\
062a175 & 2026-09-13 & Refresh full-SPD control manifest\\
8904a11 & 2026-09-13 & Record Fokker-Planck visualization utility\\
5fd0170 & 2026-09-13 & Clarify Lorenz visualization context\\
cc5eb35 & 2026-09-13 & Promote solver after full-SPD hierarchy\\
5742f7c & 2026-09-13 & Refresh full-SPD hierarchy manifest\\
e81e4fd & 2026-09-13 & Make mature bimodal certification executable\\
2cc76df & 2026-09-13 & Refresh mature certification manifest\\
d1328d7 & 2026-09-13 & Archive failed mature bimodal gate\\
0c9c51d & 2026-09-13 & Refresh mature failure manifest\\
e9d3ba9 & 2026-09-13 & Add mature positivity certificate diagnostic\\
ec11957 & 2026-09-13 & Refresh positivity diagnostic manifest\\
a3b4af5 & 2026-09-13 & Archive mature certificate diagnosis\\
e238fdb & 2026-09-13 & Refresh mature diagnosis manifest\\
592a595 & 2026-09-13 & Add mature AFC update comparator\\
ad1deef & 2026-09-14 & Add mature Bernstein DOF convex diagnostic\\
c0e413b & 2026-09-14 & Predeclare uniform mature 60 mesh reference\\
c04dc74 & 2026-09-14 & Add graded mature-state efficiency experiment\\
c91a57e & 2026-09-15 & Advance mature state and launch finer graded gate\\
eebc8b0 & 2026-09-15 & Support reliable experiment postprocess recovery\\
203186d & 2026-09-15 & Launch time-aggregated mature refinement gate\\
3aa1136 & 2026-09-15 & Start gated MFEM uniform equivalence branch\\
6bc17db & 2026-09-15 & Ignore local MFEM build products\\
9a42c4c & 2026-09-15 & Normalize MFEM conservation equivalence gate\\
bee0da6 & 2026-09-15 & Add explicit MFEM Crank Nicolson gate\\
cefdc4a & 2026-09-15 & Use ILU for nonsymmetric MFEM CN solve\\
c7c2fcd & 2026-09-15 & Calibrate MFEM CN residual tolerance\\
0f0e733 & 2026-09-15 & Gate MFEM physical equivalence before AMR\\
e6b4144 & 2026-09-15 & Match adaptive positivity certificate in MFEM gate\\
03db8af & 2026-09-16 & Add and audit MFEM mature trajectory equivalence\\
26d94f8 & 2026-09-16 & Implement gated MFEM static NC-hex mature reproduction\\
bd555b8 & 2026-09-23 & Add second-level static NC-hex mature convergence run\\
531c4e2 & 2026-09-23 & Add third-depth static AMR voxel contraction gate\\
f3f1123 & 2026-10-01 & Diagnose missed AMR support and add density-driven support test\\
ee1332a & 2026-10-01 & Add bounded depth test on fixed expanded AMR support\\
0603dce & 2026-10-02 & Visualize archived initial and evolving Lorenz densities\\
b49a35f & 2026-10-02 & Predeclare broader mature-density support sensitivity at fixed depth\\
9b92ecf & 2026-10-02 & Freeze broader mesh for mature half-timestep sensitivity\\
6255f87 & 2026-10-03 & Predeclare aligned mature-density domain sensitivity\\
84bdb48 & 2026-10-03 & Release redundant assembly storage for domain sensitivity restart\\
7e9141a & 2026-10-03 & Automatically pause and resume MPI jobs under RAM pressure\\
37f6704 & 2026-10-03 & Reuse completed first domain when restarting the second expansion\\
a1df9e8 & 2026-10-03 & Include separately grouped MPI workers in automatic RAM pausing\\
1781b40 & 2026-10-03 & Set automatic RAM resume threshold to user-requested 2 GiB\\
\bottomrule\end{longtable}\normalsize

\chapter{Source ledger and evidence index}
\label{app:ledger}
This index inventories the entire tracked tree and available compact run records. Inventory does not imply that every recorded claim is correct. Primary interpreted evidence is cited in the main chapters. The companion JSON retains SHA256 hashes, sizes, decision fields and exact commit chronology.
\section{Tracked source and preserved compact records}
\small
\begin{longtable}{p{.70\textwidth}r}
\toprule Source path & Bytes\\\midrule\endhead
\code{.gitignore} & 470\\
\code{.python-version} & 9\\
\code{AGENTS.md} & 6944\\
\code{METHOD_SELECTION_REPORT.md} & 79221\\
\code{NUMERICAL_METHOD.md} & 8022\\
\code{POSITIVITY_AUDIT.md} & 9733\\
\code{PRODUCTION_READINESS.md} & 15164\\
\code{README.md} & 10445\\
\code{benchmark_scaling.py} & 2556\\
\code{boundary_flux_report.json} & 12404\\
\code{build_experiment_manifest.py} & 5725\\
\code{corrected_q2_crank_nicolson_report.json} & 2197\\
\code{crank_nicolson_mc_20x24x24.json} & 13763\\
\code{docs/ACTIVE_TASK.md} & 143\\
\code{docs/ARCHITECTURE.md} & 5077\\
\code{docs/CLAIMS.md} & 12335\\
\code{docs/DECISIONS.md} & 40158\\
\code{docs/DENSITY_VISUALIZATION.md} & 2332\\
\code{docs/MATHEMATICAL_SPEC.md} & 4507\\
\code{docs/MATH_PROTOCOL.md} & 1895\\
\code{docs/PROJECT_STATE.md} & 16872\\
\code{docs/REFERENCES.md} & 3057\\
\code{docs/RESEARCH_DIRECTIONS.md} & 9531\\
\code{docs/assumptions.yaml} & 2245\\
\code{docs/generated/REPO_MAP.md} & 25417\\
\code{docs/tasks/TASK_TEMPLATE.md} & 1179\\
\code{environment-linux-64.lock} & 57257\\
\code{environment.yml} & 561\\
\code{experiment_manifest.json} & 25651\\
\code{experiments/corrected-q2-crank-nicolson.yaml} & 1211\\
\code{experiments/local-q2-cn-full-spd-controlled.yaml} & 2384\\
\code{experiments/local-q2-cn-full-spd-spatial-refinement.yaml} & 3429\\
\code{experiments/local-q2-cn-spatial-refinement.yaml} & 3162\\
\code{experiments/local-q2-cn-tight-ksp-refinement.yaml} & 3322\\
\code{experiments/local-q2-crank-nicolson.yaml} & 2757\\
\code{experiments/local-q2-optimizer-validation.yaml} & 1019\\
\code{experiments/local-q2-osqp-profile.yaml} & 839\\
\code{experiments/local-q2-positivity-projection.yaml} & 1805\\
\code{experiments/mature-afc-update.yaml} & 5681\\
\code{experiments/mature-dof-convex-diagnostic.yaml} & 2118\\
\code{experiments/mature-graded-fine-axes.json} & 1503\\
\code{experiments/mature-graded-fine-comparison.yaml} & 4515\\
\code{experiments/mature-graded-static-axes.json} & 1589\\
\code{experiments/mature-graded-static-comparison.yaml} & 4080\\
\code{experiments/mature-positivity-certificate-diagnostic.yaml} & 2695\\
\code{experiments/mature-state-decision.yaml} & 4573\\
\code{experiments/mature-time-aggregated-level3-pipeline.yaml} & 2588\\
\code{experiments/mature-uniform60-reference.yaml} & 2174\\
\code{experiments/method-selection-research.yaml} & 2368\\
\code{experiments/mfem-broader-support.yaml} & 683\\
\code{experiments/mfem-domain-sensitivity.yaml} & 1153\\
\code{experiments/mfem-mature-half-timestep.yaml} & 569\\
\code{experiments/mfem-mature-uniform-equivalence.yaml} & 973\\
\code{experiments/mfem-physical-equivalence-gate.yaml} & 1627\\
\code{experiments/mfem-static-nc-hex-level2.yaml} & 1622\\
\code{experiments/mfem-static-nc-hex-level3.yaml} & 1877\\
\code{experiments/mfem-static-nc-hex-reproduction.yaml} & 1728\\
\code{experiments/mfem-support-depth.yaml} & 605\\
\code{experiments/mfem-support-sensitivity.yaml} & 1267\\
\code{experiments/mfem-uniform-equivalence-gate.yaml} & 1738\\
\code{experiments/same-mesh-q2-corrected.yaml} & 1629\\
\code{experiments/same-mesh-q2-falsification.yaml} & 1235\\
\code{experiments/same-mesh-q2-timestep.yaml} & 677\\
\code{experiments/smoke-forecast.yaml} & 363\\
\code{experiments/unlimited-q2-crank-nicolson.yaml} & 1251\\
\code{experiments/unlimited-q2-timestep-attribution.yaml} & 1168\\
\code{generate_dataset.py} & 2060\\
\code{high_resolution_monte_carlo_report.json} & 3264\\
\code{higher_order_report.json} & 49522\\
\code{independent_reference_refined_report.json} & 22627\\
\code{independent_reference_report.json} & 21418\\
\code{independent_reference_study.py} & 938\\
\code{jupytext.toml} & 109\\
\code{limiter_impact_20x24x24.json} & 60388\\
\code{local_projection_optimizer_study.py} & 9438\\
\code{local_projection_study.py} & 9008\\
\code{local_q2_dynamic_projection_report.json} & 3331\\
\code{local_q2_full_spd_control_report.json} & 1691\\
\code{local_q2_full_spd_spatial_report.json} & 4246\\
\code{local_q2_mature_bimodal_report.json} & 5098\\
\code{local_q2_optimizer_validation_report.json} & 1467\\
\code{local_q2_positivity_projection_report.json} & 3880\\
\code{local_q2_projection_performance_report.json} & 3322\\
\code{local_q2_spatial_refinement_report.json} & 3215\\
\code{local_q2_tight_refinement_report.json} & 1982\\
\code{lorenz_fpe/__init__.py} & 701\\
\code{lorenz_fpe/afc.py} & 16273\\
\code{lorenz_fpe/core.py} & 70134\\
\code{lorenz_fpe/dataset.py} & 11139\\
\code{lorenz_fpe/finite_volume.py} & 6384\\
\code{lorenz_fpe/local_projection.py} & 26722\\
\code{lorenz_fpe/validation.py} & 47176\\
\code{mature_da_projected_report.json} & 440903\\
\code{mature_da_report.json} & 417269\\
\code{mature_da_study.py} & 8860\\
\code{mature_dof_convex_diagnostic.py} & 10036\\
\code{mature_forecast_validation.json} & 28581\\
\code{mature_forecast_validation.py} & 4285\\
\code{mature_graded_efficiency_report.json} & 2890\\
\code{mature_graded_fine_report.json} & 1825\\
\code{mature_positivity_certificate_diagnostic_report.json} & 4695\\
\code{mature_positivity_diagnostic.py} & 8700\\
\code{mature_time_aggregated_pipeline.py} & 13649\\
\code{mc_timestep_repair.py} & 352\\
\code{mc_timestep_same_initial_law.json} & 1795\\
\code{method_selection_report.json} & 33968\\
\code{mfem/compare_mature_equivalence.py} & 5078\\
\code{mfem/compare_physical_equivalence.py} & 4637\\
\code{mfem/compare_static_amr.py} & 6193\\
\code{mfem/compare_static_amr_level2.py} & 6016\\
\code{mfem/compare_static_amr_level3.py} & 7584\\
\code{mfem/compare_support_depth.py} & 2046\\
\code{mfem/compare_support_test.py} & 2812\\
\code{mfem/dolfinx_physical_equivalence.py} & 8098\\
\code{mfem/export_mature_initial.py} & 3617\\
\code{mfem/launch_uniform_equivalence_gate.sh} & 1171\\
\code{mfem/memory_backpressure.py} & 2647\\
\code{mfem/mfem_physical_equivalence.cpp} & 75955\\
\code{mfem/mfem_uniform_equivalence_probe.cpp} & 10665\\
\code{mfem/prepare_static_amr.py} & 3066\\
\code{mfem/prepare_static_amr_level2.py} & 3389\\
\code{mfem/prepare_static_amr_level3.py} & 3248\\
\code{mfem/prepare_support_depth.py} & 2104\\
\code{mfem/prepare_support_test.py} & 3711\\
\code{mfem/run_broader_support.py} & 5464\\
\code{mfem/run_domain_sensitivity.py} & 8061\\
\code{mfem/run_mature_uniform_equivalence.sh} & 2364\\
\code{mfem/run_physical_equivalence_gate.sh} & 1414\\
\code{mfem/run_static_amr_level2.sh} & 2869\\
\code{mfem/run_static_amr_level3.sh} & 3026\\
\code{mfem/run_static_amr_reproduction.sh} & 2768\\
\code{mfem/run_support_depth.sh} & 2035\\
\code{mfem/run_support_test.sh} & 2169\\
\code{mfem/run_uniform_equivalence_probe.sh} & 576\\
\code{monte_carlo_fixed_dt_20x24x24.json} & 13811\\
\code{monte_carlo_uncertainty_20x24x24.json} & 12815\\
\code{monte_carlo_uncertainty_30x36x36.json} & 13931\\
\code{pyproject.toml} & 545\\
\code{q1_fine_bayesian_update_report.json} & 9824\\
\code{q2_bayesian_update_report.json} & 18427\\
\code{q2_bayesian_update_subcell_report.json} & 9776\\
\code{q2_common_law_monte_carlo.json} & 14748\\
\code{q2_common_law_subcell_monte_carlo.json} & 14650\\
\code{resplit_dataset.py} & 718\\
\code{same_mesh_q2_corrected_report.json} & 2203\\
\code{same_mesh_q2_falsification_report.json} & 4473\\
\code{same_mesh_q2_study.py} & 42675\\
\code{sample_dataset/dataset_manifest.json} & 6827\\
\code{sample_dataset/dataset_schema.json} & 445\\
\code{sample_dataset/posterior_distribution_report.json} & 1592\\
\code{sample_dataset/solver_metadata.json} & 1940\\
\code{sample_dataset/trajectories/trajectory_000000/cycle_000000/diagnostics.json} & 6244\\
\code{sample_dataset/trajectories/trajectory_000000/cycle_000001/diagnostics.json} & 6748\\
\code{sample_dataset/trajectories/trajectory_000000/cycle_000002/diagnostics.json} & 6746\\
\code{sample_dataset/trajectories/trajectory_000000/cycle_000003/diagnostics.json} & 6717\\
\code{sample_dataset/trajectories/trajectory_000000/cycle_000004/diagnostics.json} & 6702\\
\code{sample_dataset/trajectories/trajectory_000000/cycle_000005/diagnostics.json} & 6723\\
\code{sample_dataset/trajectories/trajectory_000000/cycle_000006/diagnostics.json} & 6719\\
\code{sample_dataset/trajectories/trajectory_000000/cycle_000007/diagnostics.json} & 6709\\
\code{sample_dataset/trajectories/trajectory_000000/trajectory_metadata.json} & 633\\
\code{sample_dataset/trajectories/trajectory_000001/cycle_000000/diagnostics.json} & 6793\\
\code{sample_dataset/trajectories/trajectory_000001/cycle_000001/diagnostics.json} & 6589\\
\code{sample_dataset/trajectories/trajectory_000001/cycle_000002/diagnostics.json} & 6568\\
\code{sample_dataset/trajectories/trajectory_000001/cycle_000003/diagnostics.json} & 6688\\
\code{sample_dataset/trajectories/trajectory_000001/cycle_000004/diagnostics.json} & 6749\\
\code{sample_dataset/trajectories/trajectory_000001/cycle_000005/diagnostics.json} & 6759\\
\code{sample_dataset/trajectories/trajectory_000001/cycle_000006/diagnostics.json} & 6740\\
\code{sample_dataset/trajectories/trajectory_000001/cycle_000007/diagnostics.json} & 6731\\
\code{sample_dataset/trajectories/trajectory_000001/trajectory_metadata.json} & 630\\
\code{sample_dataset/trajectories/trajectory_000002/cycle_000000/diagnostics.json} & 6778\\
\code{sample_dataset/trajectories/trajectory_000002/cycle_000001/diagnostics.json} & 6753\\
\code{sample_dataset/trajectories/trajectory_000002/cycle_000002/diagnostics.json} & 6726\\
\code{sample_dataset/trajectories/trajectory_000002/cycle_000003/diagnostics.json} & 6706\\
\code{sample_dataset/trajectories/trajectory_000002/cycle_000004/diagnostics.json} & 6720\\
\code{sample_dataset/trajectories/trajectory_000002/cycle_000005/diagnostics.json} & 6725\\
\code{sample_dataset/trajectories/trajectory_000002/cycle_000006/diagnostics.json} & 6755\\
\code{sample_dataset/trajectories/trajectory_000002/cycle_000007/diagnostics.json} & 6710\\
\code{sample_dataset/trajectories/trajectory_000002/trajectory_metadata.json} & 633\\
\code{scaling_30_rank1.json} & 781\\
\code{scaling_30_rank8.json} & 778\\
\code{scaling_q2_20_rank1.json} & 793\\
\code{scaling_q2_20_rank8.json} & 795\\
\code{scaling_rank1.json} & 785\\
\code{scaling_rank2.json} & 784\\
\code{scaling_rank4.json} & 785\\
\code{scaling_rank8.json} & 790\\
\code{scripts/bootstrap} & 509\\
\code{scripts/check} & 393\\
\code{scripts/context} & 1068\\
\code{scripts/plot_saved_density.py} & 5330\\
\code{scripts/repo-map} & 606\\
\code{scripts/run-experiment} & 225\\
\code{scripts/run-in-env} & 263\\
\code{scripts/run_experiment.py} & 50153\\
\code{scripts/verify-math} & 320\\
\code{scripts/verify_math.py} & 23634\\
\code{separate_convergence_report.json} & 4004\\
\code{solver.py} & 1844\\
\code{tests/__init__.py} & 18\\
\code{tests/mpi_probe_worker.py} & 455\\
\code{tests/test_accuracy_repairs.py} & 18421\\
\code{tests/test_limiter.py} & 1501\\
\code{tests/test_memory_backpressure.py} & 1059\\
\code{tests/test_mfem_mature_compare.py} & 512\\
\code{tests/test_mpi_consistency.py} & 1265\\
\code{tests/test_solver_smoke.py} & 3761\\
\code{tests/test_spatial_experiment.py} & 6001\\
\code{tests/test_validation.py} & 834\\
\code{unlimited_q2_crank_nicolson_report.json} & 1736\\
\code{unlimited_q2_timestep_report.json} & 2350\\
\code{validate.py} & 299\\
\code{validation_repair.py} & 1239\\
\code{validation_repair_report.json} & 123701\\
\code{validation_report.json} & 81824\\
\code{voxel_projection_20x24x24.json} & 5486\\
\bottomrule\end{longtable}\normalsize
\section{Run records inspected}
The following distinct run directories have compact records in the ledger. Failed, partial and superseded runs remain visible.
\begin{itemize}
\item \code{runs/corrected-q2-crank-nicolson/20260911T002101Z}
\item \code{runs/density-visualization/20261002-saved-density}
\item \code{runs/local-q2-cn-full-spd-controlled/20260913T133432Z}
\item \code{runs/local-q2-cn-full-spd-mature-bimodal/20260913T185157Z}
\item \code{runs/local-q2-cn-full-spd-spatial-refinement/20260913T135455Z}
\item \code{runs/local-q2-cn-spatial-refinement/20260912T185711Z}
\item \code{runs/local-q2-cn-spatial-refinement/20260912T194930Z}
\item \code{runs/local-q2-cn-tight-ksp-refinement/20260912T001343Z}
\item \code{runs/local-q2-cn-tight-ksp-refinement/20260912T003213Z}
\item \code{runs/local-q2-crank-nicolson/20260911T115805Z}
\item \code{runs/local-q2-crank-nicolson/20260911T165516Z}
\item \code{runs/local-q2-optimizer-validation/20260911T165132Z}
\item \code{runs/local-q2-optimizer-validation/20260911T165327Z}
\item \code{runs/local-q2-osqp-profile/20260911T165425Z}
\item \code{runs/local-q2-positivity-projection/20260911T114352Z}
\item \code{runs/local-q2-positivity-projection/20260911T114732Z}
\item \code{runs/local-q2-positivity-projection/20260911T114957Z}
\item \code{runs/local-q2-spatial-runner-smoke/20260912T185453Z}
\item \code{runs/mature-dof-convex-diagnostic/20260914T184021Z}
\item \code{runs/mature-full-spd-afc-update/20260913T222328Z}
\item \code{runs/mature-full-spd-graded-fine-comparison/20260914T231639Z}
\item \code{runs/mature-full-spd-graded-static-comparison/20260914T210354Z}
\item \code{runs/mature-full-spd-uniform60-reference/20260914T190730Z}
\item \code{runs/mature-positivity-certificate-diagnostic/20260913T200251Z}
\item \code{runs/mature-time-aggregated-grid-pipeline/20260915T134916Z}
\item \code{runs/mfem-broader-support/20261002-preflight}
\item \code{runs/mfem-broader-support/20261002T080644Z}
\item \code{runs/mfem-domain-memory-regression/20261003-after}
\item \code{runs/mfem-domain-memory-regression/20261003-before}
\item \code{runs/mfem-domain-sensitivity/20261003T105837Z}
\item \code{runs/mfem-domain-sensitivity/20261003T113423Z}
\item \code{runs/mfem-domain-sensitivity/20261003T143701Z}
\item \code{runs/mfem-domain-sensitivity/20261003T143753Z}
\item \code{runs/mfem-mature-half-timestep/20261002-preflight}
\item \code{runs/mfem-mature-half-timestep/20261002T171744Z}
\item \code{runs/mfem-mature-uniform-equivalence/20260915T192747Z}
\item \code{runs/mfem-mature-uniform-equivalence/20260915T192800Z}
\item \code{runs/mfem-mature-uniform-equivalence/20260915T192954Z}
\item \code{runs/mfem-mature-uniform-equivalence/20260915T193057Z}
\item \code{runs/mfem-mature-uniform-equivalence/20260915T193238Z}
\item \code{runs/mfem-mature-uniform-equivalence/20260915T193445Z}
\item \code{runs/mfem-mature-uniform-equivalence/20260916T213458Z}
\item \code{runs/mfem-nc-face-probe/20260916T-probe}
\item \code{runs/mfem-physical-equivalence-production/20260915T184539Z}
\item \code{runs/mfem-physical-equivalence-production/20260915T184604Z}
\item \code{runs/mfem-physical-equivalence-production/20260915T185006Z}
\item \code{runs/mfem-physical-equivalence/20260915T182726Z}
\item \code{runs/mfem-physical-equivalence/20260915T182747Z}
\item \code{runs/mfem-physical-equivalence/20260915T182953Z}
\item \code{runs/mfem-physical-equivalence/20260915T183056Z}
\item \code{runs/mfem-physical-equivalence/20260915T183234Z}
\item \code{runs/mfem-physical-equivalence/20260915T183333Z}
\item \code{runs/mfem-physical-equivalence/20260915T183412Z}
\item \code{runs/mfem-physical-equivalence/20260915T183510Z}
\item \code{runs/mfem-physical-equivalence/20260915T183709Z}
\item \code{runs/mfem-physical-equivalence/20260915T184057Z}
\item \code{runs/mfem-physical-equivalence/20260915T184304Z}
\item \code{runs/mfem-physical-equivalence/20260915T184941Z}
\item \code{runs/mfem-physical-equivalence/20260915T184954Z}
\item \code{runs/mfem-physical-equivalence/20260915T193047Z}
\item \code{runs/mfem-physical-equivalence/20260915T193227Z}
\item \code{runs/mfem-physical-equivalence/20260916T213441Z}
\item \code{runs/mfem-physical-equivalence/20260916T220950Z}
\item \code{runs/mfem-physical-equivalence/serial-debug-20260915T183529Z}
\item \code{runs/mfem-static-amr-level2/20260922T231219Z}
\item \code{runs/mfem-static-amr-level3/20260923T181337Z}
\item \code{runs/mfem-static-amr-reproduction/20260916T221235Z}
\item \code{runs/mfem-static-amr-reproduction/design-20260916}
\item \code{runs/mfem-static-amr-smoke/20260916}
\item \code{runs/mfem-static-amr-smoke/20260916-amr45}
\item \code{runs/mfem-static-amr-smoke/20260916-coverage}
\item \code{runs/mfem-static-amr-smoke/20260916-uniform30}
\item \code{runs/mfem-support-depth/20261001T224713Z}
\item \code{runs/mfem-support-sensitivity/20261001T144235Z}
\item \code{runs/mfem-uniform-equivalence-gate/20260915T163908Z}
\item \code{runs/mfem-uniform-equivalence-gate/20260915T165354Z}
\item \code{runs/mfem-uniform-equivalence-gate/20260915T170152Z}
\item \code{runs/mfem-uniform-equivalence-gate/20260915T170429Z}
\item \code{runs/same-mesh-q2-corrected/20260910T234135Z}
\item \code{runs/same-mesh-q2-falsification/20260910T185420Z}
\item \code{runs/same-mesh-q2-recovery/20260910T190900Z}
\item \code{runs/same-mesh-q2-timestep/20260910T190840Z}
\item \code{runs/smoke-forecast-q1/20260909T020228Z}
\item \code{runs/smoke-forecast-q1/20260910T185328Z}
\item \code{runs/unlimited-q2-crank-nicolson/20260911T001508Z}
\item \code{runs/unlimited-q2-timestep-attribution/20260911T000813Z}
\end{itemize}
\section{Independent recomputation}
\begin{lstlisting}[basicstyle=\ttfamily\footnotesize]
{
  "l1": 1.9549638427469373e-12,
  "mass": 0.9999999999986662,
  "negative_mass": 0.0,
  "added_shell_mass": 0.0,
  "finite": true,
  "input_hashes": {
    "runs/mfem-domain-sensitivity/20261003T143753Z/padding2/mfem/final_q2_subcell_averages.bin": "2014ef4a230cb172ff3235b1df243927e41fc5eb42e8a41b0fd18200d1a0a455",
    "runs/mfem-domain-sensitivity/20261003T113423Z/padding1/mfem/final_q2_subcell_averages.bin": "af7d4d34f3fca596c6f9ec4e2ea02485a3de37ee51739c8dcb7ee2655ebe8875"
  }
}
\end{lstlisting}

\backmatter
\begin{thebibliography}{99}
\addcontentsline{toc}{chapter}{Bibliography}
\bibitem{lorenz} E. N. Lorenz. Deterministic nonperiodic flow. \emph{Journal of the Atmospheric Sciences} 20(2), 130--141, 1963. DOI: \href{https://doi.org/10.1175/1520-0469(1963)020%3C0130:DNF%3E2.0.CO;2}{10.1175/1520-0469(1963)020\textless0130:DNF\textgreater2.0.CO;2}. Publisher record: \url{https://journals.ametsoc.org/doi/abs/10.1175/1520-0469%281963%29020%3C0130%3Adnf%3E2.0.CO%3B2}.
\bibitem{sarkka} S. S\"arkk\"a and A. Solin. \emph{Applied Stochastic Differential Equations}. Cambridge University Press, 2019. Author-hosted book: \url{https://users.aalto.fi/~asolin/sde-book/sde-book.pdf}. Used for the standard SDE/generator context; the project-specific derivations are given in this monograph.
\bibitem{chang} J. S. Chang and G. Cooper. A practical difference scheme for Fokker--Planck equations. \emph{Journal of Computational Physics} 6(1), 1--16, 1970. DOI: \href{https://doi.org/10.1016/0021-9991(70)90001-X}{10.1016/0021-9991(70)90001-X}. This is literature context, not evidence that a Chang--Cooper branch was run in the present repository.
\bibitem{arnold} D. N. Arnold, F. Brezzi, B. Cockburn and L. D. Marini. Unified analysis of discontinuous Galerkin methods for elliptic problems. \emph{SIAM Journal on Numerical Analysis} 39(5), 1749--1779, 2002. DOI: \href{https://doi.org/10.1137/S0036142901384162}{10.1137/S0036142901384162}. Applies with appropriate coercivity, mesh and regularity hypotheses; it is not a theorem for the complete nonlinear corrected forecast.
\bibitem{kuzmin} D. Kuzmin. On the design of general-purpose flux limiters for finite element schemes. I. Scalar convection. \emph{Journal of Computational Physics} 219(2), 513--531, 2006. DOI: \href{https://doi.org/10.1016/j.jcp.2006.03.034}{10.1016/j.jcp.2006.03.034}. General AFC context; the project comparator's complete-graph average correction is identified separately.
\bibitem{carrillo} J. A. Carrillo, H. Liu and H. Yu. Positivity-preserving and energy-dissipating discontinuous Galerkin methods for nonlinear nonlocal Fokker--Planck equations. \emph{Communications in Applied and Industrial Mathematics} 16(1), 19--40, 2025. DOI: \href{https://doi.org/10.2478/caim-2025-0002}{10.2478/caim-2025-0002}. Author preprint: \url{https://arxiv.org/abs/2403.15643}. Their energy structure and DG construction differ from the Lorenz SIPG/CN branch.
\bibitem{liuhu} C. Liu, J. Hu, W. T. Taitano and X. Zhang. An optimization-based positivity-preserving limiter in semi-implicit discontinuous Galerkin schemes solving Fokker--Planck equations. \emph{Computers \& Mathematics with Applications} 192, 54--71, 2025. DOI: \href{https://doi.org/10.1016/j.camwa.2025.05.008}{10.1016/j.camwa.2025.05.008}. Preprint: \url{https://arxiv.org/abs/2410.19143}. A neighbouring optimisation architecture; their temporal/flux construction is not identical to the current method.
\bibitem{osqp} B. Stellato, G. Banjac, P. Goulart, A. Bemporad and S. Boyd. OSQP: an operator splitting solver for quadratic programs. \emph{Mathematical Programming Computation} 12(4), 637--672, 2020. DOI: \href{https://doi.org/10.1007/s12532-020-00179-2}{10.1007/s12532-020-00179-2}. Author record: \url{https://web.stanford.edu/~boyd/papers/osqp.html}. Software documentation: \url{https://osqp.org/docs/}.
\bibitem{mfemnc} MFEM project. Nonconforming adaptive mesh refinement. Official documentation: \url{https://mfem.org/howto/ncmesh/}. Consulted 3 October 2026. Supports the infrastructure description; numerical correctness of the project fluxes is established by its separate probes.
\bibitem{mfemdiff} MFEM project. \code{DGDiffusionIntegrator}, version 4.8 code documentation. \url{https://docs.mfem.org/4.8/classmfem_1_1DGDiffusionIntegrator.html}. Matrix-coefficient SIPG interface; consult the archived project build for exact ABI and implementation settings.
\bibitem{linuxproc} Linux kernel project. The /proc filesystem: memory information and \code{MemAvailable}. \url{https://docs.kernel.org/filesystems/proc.html}. Consulted 3 October 2026. Kernel memory estimates do not constitute a solver checkpoint guarantee.
\bibitem{project} A. J. Sheppard's Lorenz Fokker--Planck research repository. Solver revision \code{1781b40ad156f01702aeeb16eaf7e0240ede7378}, 3 October 2026. \url{https://github.com/AdamJamesSheppard/lorenz-training-set}. Primary local numerical evidence is indexed in \code{source_ledger.json}; ignored arrays are not all distributed by the Git repository.
\end{thebibliography}

\end{document}
